% !TEX root = ../main.tex
\chapter{Stochastic Processes II: Brownian Motion}\label{ch:sp2}
\begin{paracol}{2}

\begin{sources}
\begin{tabularx}{\linewidth}{@{}l l L@{}}
\toprule
\textbf{Lecture} & \textbf{Speaker} & \textbf{Content}\\
\midrule
2024 L14 & P.~Kempthorne & History; Brownian motion with diffusion coefficient $\sigma^2$: increments, Markov property, covariance; transition density and the diffusion equation; Brownian motion as limit of normalised random walks (invariance principle); reflection principle, maximum, first hitting time; reflected and absorbed Brownian motion; Brownian bridge and the empirical distribution function; Brownian motion with drift; gambler's ruin by infinitesimal one-step analysis; non-differentiability; quadratic variation; log-price quadratic variation; Laplace increments at exponential times; filtrations and adapted processes.\\
2013 L17 & C.~Lee & Continuous-time processes and why their law is hard to describe; the existence theorem for standard Brownian motion; random-walk limit; Brown, Einstein and Bachelier; zero crossings, $\sqrt t$ envelope, nowhere differentiability; distribution of the maximum by the reflection principle; quadratic variation and $(\dd B)^2=\dd t$; Brownian motion with drift; \emph{why $S_t\ne e^{B_t}$}, Taylor expansion and the first appearance of It\^o's lemma.\\
\bottomrule
\end{tabularx}

\smallskip
\textbf{Files.} 2024 slides \file{mit18_642_f24_lec14_1.pdf} (77 pages with overlays, 20 slides) and the figure decks
\file{lec14_2} (two-dimensional Brownian motion, 201 animation frames), \file{lec14_3} (normal densities $p(x\mid t)$,
$t=1,\dots,10$), \file{lec14_4} (limiting random walk, 201 frames), \file{lec14_5} (reflected Brownian motion, three trials),
\file{lec14_6} ($\Prob[M(t)>x]$), \file{lec14_7} (hitting-time density), \file{lec14_8} (gambler's ruin with drift, four
panels), \file{lec14_9} (quadratic variation, 201 frames); 2013 note \file{MIT18_S096F13_lecnote17.pdf}.
\textbf{Problem sets.} 2024 PS4 Problems~1--2 (conditional moments and scaling of Brownian motion; Problems~3--4 are the
maximum-likelihood estimation of a geometric Brownian motion, \cref{ch:volatility}); 2024 PS5 Problems~1--2 (L\'evy processes:
Brownian motion with drift and the Poisson process, with simulations; Problems~3--5 are volatility estimation,
\cref{ch:volatility}). 2013 PS8 Problem~A-2 is 2024 PS4 Problem~1 and Problem~B-1 is on two-dimensional Brownian motion; PS8
belongs to \cref{ch:stochcalc}, see the course note in \cref{ch11:sec-ahead}. The exercises are at the end of the chapter.
\textbf{Not covered.} The construction of Brownian motion (the existence theorem is only quoted in both years), the
proofs of Donsker's theorem and of nowhere differentiability, and the law of the iterated logarithm.
The R code that generated the figure decks is not part of the download; the figures of this chapter are recomputed from formulas.
\end{sources}

\section*{Motivation}
Chapter~\ref{ch:sp1} studied random walks, martingales and stopping times in discrete time. Prices, however,
change at every instant a trade takes place, and the mathematics of derivative pricing (\cref{ch:stochcalc,ch:sde,ch:bs}) is
carried by a \emph{continuous-time} model. Its basic building block is \emph{Brownian motion}, the continuous-time random walk.
It was found three times, in three unrelated places: by the botanist Robert Brown (1827), watching pollen grains jitter in
water; by Louis Bachelier (1900), as a model of prices on the Paris exchange; by Albert Einstein (1905), who explained
Brown's observation by the bombardment of the grain by water molecules and obtained the diffusion law
$\langle x^2\rangle=2Dt$ \ts{24}{14}{0:00}\ts{13}{17}{18:21}\ts{13}{17}{20:29}. Norbert Wiener, an MIT mathematician,
gave it its rigorous foundation in 1923, which is why it is also called the \emph{Wiener process}
\ts{24}{14}{1:04}\ts{13}{17}{13:43}.

The reason for its centrality is a \emph{universality} statement, made precise in \cref{ch11:sec-donsker}: any random walk with
finite-variance steps, watched from far away, looks like Brownian motion. For a physicist this is the central limit theorem in
disguise; for a trader it says that a price built from many small, roughly independent trades is described by two numbers, a
drift and a volatility. This chapter develops the calculus of \emph{sample paths} that the next chapters use:

\begin{itemize}
  \item the definition and the elementary properties (Gaussian, Markov, martingale), and the diffusion equation
        (\cref{ch11:sec-def});
  \item Brownian motion as the scaling limit of the random walk (\cref{ch11:sec-donsker});
  \item the strange geometry of the paths: continuous but nowhere differentiable, with \emph{finite quadratic variation}
        $[B]_T=T$, the fact that later becomes It\^o's lemma (\cref{ch11:sec-paths});
  \item the reflection principle, giving the law of the running maximum and of the first hitting time
        (\cref{ch11:sec-reflection}), with reflected and absorbed Brownian motion;
  \item Brownian motion with drift and the gambler's ruin problem, geometric Brownian motion, and why the
        price is \emph{not} $e^{B_t}$ (\cref{ch11:sec-drift,ch11:sec-gbm});
  \item further processes built from $B$: bridge, random-time sampling, two dimensions, L\'evy processes, filtrations
        (\cref{ch11:sec-variants}).
\end{itemize}

% ==================================================================
\section{Definition and elementary properties}\label{ch11:sec-def}

\subsection{Brownian motion}
In discrete time the law of a random walk is specified through its increments: ``each step is $\pm1$ with probability
$\tfrac12$''. In continuous time there is no next step to describe, so one lists \emph{properties} of increments over all
intervals and asserts that a process with these properties exists \ts{13}{17}{4:10}\ts{13}{17}{5:18}\ts{13}{17}{9:28}.

\begin{definition}[Standard Brownian motion]\label{ch11:def-bm}
A stochastic process $\{B(t),\,t\ge0\}$ is a \emph{standard Brownian motion} (Wiener process) if
\begin{enumerate}[(i)]
  \item $B(0)=0$ and $t\mapsto B(t)$ is continuous (with probability one);
  \item (stationary Gaussian increments) for $0\le s<t$, $B(t)-B(s)\sim N(0,\,t-s)$;
  \item (independent increments) for $t_1\le t_2\le t_3\le t_4$ the increments $B(t_2)-B(t_1)$ and $B(t_4)-B(t_3)$ are
        independent; more generally the increments over non-overlapping intervals are mutually independent
        \ts{24}{14}{3:08}\ts{24}{14}{4:13}\ts{24}{14}{5:17}\ts{13}{17}{7:21}.
\end{enumerate}
\end{definition}

The 2024 slides call $\sigma^2$ the \emph{diffusion coefficient} and write $B(t+\Delta t)-B(t)\sim N(0,\sigma^2\Delta t)$
already in the definition; from the slide on drift onwards $B$ is standard and $\sigma B$ is used. We follow the second convention:
\emph{Brownian motion with diffusion coefficient $\sigma^2$} is $\sigma B(t)$, and one may start it at any level $x$ by
adding $x$ (``Brownian motion started at $x$'') \ts{24}{14}{9:36}. Property (ii) is the one to remember: over a time span
$h$ the process moves by a centred normal variable whose \emph{variance is the time elapsed}, so its typical size is
$\sqrt h$ \ts{13}{17}{12:40}.

\begin{coursenote}[Existence, and why the definition is indirect]
Both lecturers stress that (i)--(iii) are a \emph{list of demands} on a probability law on the space of continuous
paths, and that the theorem ``such a law exists'' is hard and not proved in the course \ts{13}{17}{6:20}\ts{13}{17}{10:36}\ts{13}{17}{11:37}.
The sample space is the space of paths $\omega:[0,\infty)\to\R$ itself, with $\Prob$ a measure on it (Wiener measure); a
path is an outcome, and there is no ``density at $\omega$'' \ts{13}{17}{5:18}. One standard route (Wiener 1923; L\'evy 1948)
builds $B$ as a uniform limit of random piecewise-linear functions obtained by refining a dyadic grid; another is
Donsker's theorem below, where $B$ is \emph{defined} as the limit of rescaled random walks (this is how the 2024 lecture
introduces it, \ts{24}{14}{17:22}). What (i)--(iii) do determine is the joint law of $(B(t_1),\dots,B(t_m))$ for every finite
set of times, by writing $B(t_k)$ as a sum of independent normal increments; continuity of the paths is the extra input
that makes the law unique on path space.
\end{coursenote}

\begin{proposition}[Gaussian characterisation]\label{ch11:prop-gauss}
A process with continuous paths and $B(0)=0$ is a standard Brownian motion if and only if it is a centred Gaussian process
with covariance function $\Cov(B(s),B(t))=\min(s,t)$.
\end{proposition}
\begin{proof}
($\Rightarrow$) Take $s\le t$. Then $B(t)=B(s)+[B(t)-B(s)]$ with independent summands, so
$\E[B(s)B(t)]=\E[B(s)^2]+\E[B(s)]\,\E[B(t)-B(s)]=s+0$ \ts{24}{14}{7:33}\ts{24}{14}{8:34}. Any vector $(B(t_1),\dots,B(t_m))$ is a linear
image of independent normal increments, hence jointly Gaussian.
($\Leftarrow$) For $s<t$ the variable $B(t)-B(s)$ is centred normal with variance $t-2s+s=t-s$. For
$s_1<t_1\le s_2<t_2$, $\Cov(B(t_1)-B(s_1),B(t_2)-B(s_2))=t_1-t_1-s_1+s_1=0$ (write out the four
minima); jointly Gaussian and uncorrelated implies independent.
\end{proof}
Note the meaning of the covariance: $B(s)$ and $B(t)$ are correlated to the extent that the two time windows $[0,s]$ and $[0,t]$
overlap, $\Corr(B(s),B(t))=\sqrt{s/t}$ for $s<t$ \ts{24}{14}{8:34}.

\begin{proposition}[Elementary properties]\label{ch11:prop-elem}
Let $B$ be a standard Brownian motion and $\mathcal F_t=\sigma(B(u),u\le t)$.
\begin{enumerate}[(i)]
  \item \emph{Mean and variance.} $\E[B(t)]=0$, $\Var B(t)=t$: the standard deviation grows like $\sqrt t$
        \ts{24}{14}{6:28}\ts{24}{14}{7:33}.
  \item \emph{Markov property.} For $h>0$, $B(t+h)=B(t)+[B(t+h)-B(t)]$ where the bracket is $N(0,h)$ and independent of
        $\mathcal F_t$; hence the law of $B(t+h)$ given $\mathcal F_t$ is $N(B(t),h)$, which depends on the past only through $B(t)$
        \ts{24}{14}{5:17}\ts{24}{14}{6:28}\ts{24}{14}{11:54}.
  \item \emph{Martingales.} $B(t)$, $B(t)^2-t$ and, for every real $\lambda$, $\exp(\lambda B(t)-\tfrac12\lambda^2t)$ are
        martingales with respect to $\{\mathcal F_t\}$.
  \item \emph{Symmetries.} $-B$; the shifted process $B(t_0+t)-B(t_0)$; the rescaled process $c^{-1/2}B(ct)$ for $c>0$; and the
        time-inverted process $tB(1/t)$ ($t>0$, and $0$ at $t=0$) are again standard Brownian motions.
\end{enumerate}
\end{proposition}
\begin{proof}
(i) is property (ii) of \cref{ch11:def-bm} with $s=0$.

(ii) is property (iii) of \cref{ch11:def-bm} together with the definition of $\mathcal F_t$.

(iii) For $s<t$, using independence of
$B(t)-B(s)$ from $\mathcal F_s$: $\E[B(t)\mid\mathcal F_s]=B(s)$;
$\E[B(t)^2\mid\mathcal F_s]=B(s)^2+2B(s)\E[B(t)-B(s)]+\E[(B(t)-B(s))^2]=B(s)^2+(t-s)$; and with the moment generating
function of a normal variable, $\E[e^{\lambda(B(t)-B(s))}]=e^{\lambda^2(t-s)/2}$, so
$\E[e^{\lambda B(t)-\lambda^2t/2}\mid\mathcal F_s]=e^{\lambda B(s)}e^{\lambda^2(t-s)/2}e^{-\lambda^2t/2}=e^{\lambda B(s)-\lambda^2s/2}$.

(iv) Each candidate is continuous (at $t=0$ for the inverted one use $B(u)/u\to0$, which is the law of large numbers for
$B$), centred Gaussian, and one checks that its covariance is $\min(s,t)$; then apply \cref{ch11:prop-gauss}.
For the rescaling, $\Cov(c^{-1/2}B(cs),c^{-1/2}B(ct))=c^{-1}\min(cs,ct)=\min(s,t)$; this is the content of 2024 PS4
Problem~2 (\cref{ch11:ex-scaling}); for the inversion, $st\min(1/s,1/t)=\min(t,s)$.
\end{proof}

Property (c) is the continuous-time image of \cref{ch04:ex-m1,ch04:ex-m2,ch04:ex-m4}: the same three martingales that
appeared for the random walk (sum, square minus $n\sigma^2$, exponential) survive the limit. The variance $t$ that
appears in $B^2-t$ is the \emph{quadratic variation} of \cref{ch11:sec-paths}. The martingale $\exp(\lambda B-\lambda^2t/2)$ is
the ancestor of the change of measure behind risk-neutral pricing (\cref{ch:bs}).

\begin{intuition}[The physicist's reading]\label{ch11:int-physics}
Standard Brownian motion is the position of a free overdamped particle whose velocity has been eliminated: the Langevin
equation $\dot x=\sqrt{2D}\,\xi(t)$ with white noise $\xi$ has solution $x(t)=\sqrt{2D}\,B(t)$, so
$\langle x^2\rangle=2Dt$ (Einstein's law) and $\sigma^2=2D$. The white noise $\xi=\dot B$ does not exist as a function
(\cref{ch11:sec-paths}); only its integral does. The law of the whole path is the \emph{Wiener measure}, the measure of the
Euclidean path integral of a free particle,
$\propto\exp\bigl(-\tfrac12\int_0^t\dot x^2\dd s\bigr)\prod\dd x(s)$, formal but useful; the transition density below is
its propagator (the heat kernel), and $\min(s,t)$ is the Green function of $-\dd^2/\dd t^2$ with Dirichlet condition at $0$.
\end{intuition}

\begin{finance}[Bachelier's model]
Bachelier's thesis modelled the price itself as $S(t)=S(0)+\mu t+\sigma B(t)$, ``arithmetic'' Brownian motion; the
2013 note credits him with the first use of the process to value stocks and options \ts{13}{17}{13:43}. Its defect is that the price becomes
negative with positive probability, so equities are modelled through $\log S$ or through relative changes
(\cref{ch11:sec-gbm}). The arithmetic model survives where levels can be negative or small compared with their moves, in
particular for interest rates and spreads, where ``normal volatility'' is quoted in basis points per year.
\end{finance}

% ------------------------------------------------------------------
\subsection{Transition density and the diffusion equation}\label{ch11:sec-diffusion}
Let $B^x$ be Brownian motion started at $x$, i.e.\ $x+B$. For $\tau>0$ and diffusion coefficient $\sigma^2$, the density of
$B^x(s+\tau)$ given $B^x(s)=x$ is the Gaussian
\begin{equation}\label{ch11:eq-kernel}
  p(y,\tau\mid x)=\frac1{\sqrt{2\pi\sigma^2\tau}}\exp\Bigl(-\frac{(y-x)^2}{2\sigma^2\tau}\Bigr),
\end{equation}
that is, $\Prob[B(t)\le c\mid B(s)=x]=\Phi\bigl((c-x)/(\sigma\sqrt{t-s})\bigr)$ in terms of the standard normal cdf $\Phi$ and pdf
$\varphi$ \ts{24}{14}{9:36}\ts{24}{14}{10:42}\ts{24}{14}{11:54}\ts{24}{14}{13:03}.

\begin{proposition}[Diffusion equation]\label{ch11:prop-heat}
The kernel \eqref{ch11:eq-kernel} satisfies
\begin{equation}\label{ch11:eq-heat}
  \frac{\partial p}{\partial\tau}=\frac{\sigma^2}2\,\frac{\partial^2p}{\partial y^2},
  \qquad \int_{-\infty}^{\infty}p(y,\tau\mid x)\dd y=1\ \ (\tau>0),\qquad p(y,\tau\mid x)\xrightarrow[\tau\downarrow0]{}\delta(y-x),
\end{equation}
and \eqref{ch11:eq-heat} has no other solution with these properties. The same holds in the variable $x$ (backward equation).
\end{proposition}
\begin{proof}
Put $u=y-x$. Then $\partial_\tau p=p\,\bigl(-\tfrac1{2\tau}+\tfrac{u^2}{2\sigma^2\tau^2}\bigr)$ and
$\partial_y^2p=p\,\bigl(-\tfrac1{\sigma^2\tau}+\tfrac{u^2}{\sigma^4\tau^2}\bigr)$, so $\partial_\tau p=\tfrac{\sigma^2}2\partial_y^2p$.
Integrating \eqref{ch11:eq-heat} over $y$ gives $\partial_\tau\int p=\tfrac{\sigma^2}2[\partial_yp]_{-\infty}^{\infty}=0$: probability
is conserved \ts{24}{14}{16:17}. The Fourier transform in $y$ turns the equation into $\partial_\tau\hat p=-\tfrac12\sigma^2k^2\hat p$,
whose solution with $\hat p(k,0)=e^{-ikx}$ is $e^{-ikx-\sigma^2k^2\tau/2}$, the transform of \eqref{ch11:eq-kernel}; this gives uniqueness
among tempered solutions \ts{24}{14}{17:22}. The backward version follows from $p$ depending on $y-x$ only.
\end{proof}
The semigroup property (Chapman--Kolmogorov, \cref{ch04:thm-nstep}) is the convolution identity
$\int p(z,\tau_2\mid y)\,p(y,\tau_1\mid x)\dd y=p(z,\tau_1+\tau_2\mid x)$, the statement that a sum of independent normals is normal
with added variances. In other disciplines \eqref{ch11:eq-heat} is the \emph{heat equation}, in population dynamics the
\emph{Fokker--Planck} or \emph{Kolmogorov forward} equation \ts{24}{14}{17:22}; it returns in the valuation of derivatives
(\cref{ch:bs}), where the Black--Scholes PDE is a backward equation of this type.

\begin{figure}[ht]
\centering
\begin{tikzpicture}
\begin{axis}[width=0.9\linewidth,height=5cm,xlabel={$y$},ylabel={$p(y,t\mid0)$},xmin=-6,xmax=6,ymin=0,ymax=0.45,
  legend style={at={(0.5,1.03)},anchor=south,legend columns=3,font=\footnotesize},grid=major,grid style={black!10}]
  \addplot[thick,color=defcol,domain=-6:6,samples=100] {1/sqrt(2*pi*1)*exp(-x^2/(2*1))};
  \addplot[thick,color=intcol,domain=-6:6,samples=100] {1/sqrt(2*pi*2.5)*exp(-x^2/(2*2.5))};
  \addplot[thick,color=thmcol,domain=-6:6,samples=100] {1/sqrt(2*pi*6)*exp(-x^2/(2*6))};
  \addplot[only marks,mark=*,mark size=1.6pt,color=defcol] coordinates {(-1,0.24197) (1,0.24197)};
  \addplot[only marks,mark=*,mark size=1.6pt,color=intcol] coordinates {(-1.5811,0.15316) (1.5811,0.15316)};
  \addplot[only marks,mark=*,mark size=1.6pt,color=thmcol] coordinates {(-2.4495,0.09946) (2.4495,0.09946)};
  \legend{{$t=1$},{$t=2.5$},{$t=6$}}
\end{axis}
\end{tikzpicture}
\caption{Diffusion of the density of $B(t)$ (the panels of \file{lec14_3}, which shows $t=1,1.5,\dots,10$). The
density drops at the centre and rises in the tails; the dots mark the inflection points $y=\pm\sqrt t$, where the curvature
changes sign, so that $\partial_tp<0$ inside and $>0$ outside \ts{24}{14}{14:08}\ts{24}{14}{15:14}.}\label{ch11:fig-heat}
\end{figure}


% ==================================================================
\section{Brownian motion as the limit of random walks}\label{ch11:sec-donsker}

Let $X_1,X_2,\dots$ be i.i.d.\ with $\E[X_i]=0$, $\Var X_i=1$, and $S_n=X_1+\dots+X_n$
(\cref{ch04:sec-rw}). By the central limit theorem (\cref{ch03:thm-clt}), $\Prob[S_n/\sqrt n\le c]\to\Phi(c)$, i.e.\ $\Prob[S_n\le x]\approx\Phi(x/\sqrt n)$
\ts{24}{14}{17:22}\ts{24}{14}{18:27}. To speak about a \emph{path} rather than one time, rescale space by $\sqrt n$ and time by $n$.

\begin{definition}[Normalised random walk]\label{ch11:def-bn}
$B_n(t)=S_{\lfloor nt\rfloor}/\sqrt n$ for $t\ge0$, where $\lfloor nt\rfloor$ is the largest integer $\le nt$; the process jumps at
$t=k/n$ by $X_k/\sqrt n$ and is constant in between \ts{24}{14}{19:29}. The piecewise-linear version, which joins the points
$(k/n,S_k/\sqrt n)$ (the version of the 2013 note, ``connect the dots''), is continuous and differs from it by at most
$\max_{k\le nT}|X_k|/\sqrt n$ on $[0,T]$ \ts{13}{17}{14:56}\ts{13}{17}{16:01}.
\end{definition}

\begin{coursenote}[Erratum in the 2013 note]
\file{lecnote17} defines the piecewise-linear path by $Z(t/n)=Y_t$ for a walk with variance-one steps; the rescaling by $1/\sqrt n$ is missing (the
values would be of order $\sqrt n$ and no limit could exist). The correct definition is $Z(t/n)=Y_t/\sqrt n$, as in the 2024 slides.
\end{coursenote}

\begin{theorem}[Donsker's invariance principle]\label{ch11:thm-donsker}
As $n\to\infty$ the process $B_n$ converges in distribution (weakly, as a random element of the space of paths on $[0,T]$
with the uniform metric) to a standard Brownian motion $B$, whatever the distribution of the steps $X_i$
\ts{24}{14}{20:44}. In particular, for every continuous functional $F$ of the path (e.g.\ $\sup_{t\le1}$, $\int_0^1(\cdot)^2\dd t$,
the value at time $1$), $F(B_n)\xrightarrow{d}F(B)$.
\end{theorem}
\begin{proof}[What is proved here]
\emph{Finite-dimensional convergence.} For $0\le t_1<\dots<t_m$ the increments $B_n(t_j)-B_n(t_{j-1})$ are sums of
$\lfloor nt_j\rfloor-\lfloor nt_{j-1}\rfloor\sim n(t_j-t_{j-1})$ i.i.d.\ steps divided by $\sqrt n$, over disjoint blocks, hence independent; by the central limit
theorem each converges to $N(0,t_j-t_{j-1})$, and by independence jointly, which is exactly the joint law of the increments of $B$
\ts{13}{17}{17:19}. This gives $\Var B(t)=\lim_n\Var B_n(t)=t$ and $\E[B(t)]=0$, the mean and variance of the slides
\ts{24}{14}{19:29}. \emph{Tightness}, that no path of $B_n$ oscillates wildly between grid times so that the limit is a
continuous process, requires a maximal inequality (Kolmogorov or Doob) and is not proved in the course; it is what turns
convergence of the finite-dimensional laws into convergence of laws on path space \ts{13}{17}{6:20}.
\end{proof}

\begin{figure}[ht]
\centering
\begin{tikzpicture}
\begin{axis}[width=0.47\linewidth,height=4.8cm,xlabel={$t$},ylabel={$B_n(t)$},xmin=0,xmax=1,title={$n=20$, $\pm1$ steps},title style={font=\footnotesize},grid=major,grid style={black!10}]
  \addplot[thick,const plot,color=defcol] coordinates {(0.000,0.000) (0.050,-0.224) (0.100,-0.447) (0.150,-0.224) (0.200,-0.447) (0.250,-0.224) (0.300,0.000) (0.350,0.224) (0.400,0.000) (0.450,-0.224) (0.500,-0.447) (0.550,-0.671) (0.600,-0.447) (0.650,-0.224) (0.700,-0.447) (0.750,-0.224) (0.800,-0.447) (0.850,-0.224) (0.900,0.000) (0.950,0.224) (1.000,0.447)};
\end{axis}
\end{tikzpicture}\hfill
\begin{tikzpicture}
\begin{axis}[width=0.47\linewidth,height=4.8cm,xlabel={$t$},xmin=0,xmax=1,title={$n=150$},title style={font=\footnotesize},
  legend style={at={(0.5,1.25)},anchor=south,legend columns=2,font=\scriptsize},grid=major,grid style={black!10}]
  \addplot[thick,color=defcol] coordinates {(0.000,0.000) (0.007,0.082) (0.013,0.163) (0.020,0.082) (0.027,0.163) (0.033,0.082) (0.040,0.163) (0.047,0.245) (0.053,0.163) (0.060,0.245) (0.067,0.163) (0.073,0.082) (0.080,0.000) (0.087,0.082) (0.093,0.000) (0.100,-0.082) (0.107,-0.163) (0.113,-0.245) (0.120,-0.327) (0.127,-0.408) (0.133,-0.327) (0.140,-0.408) (0.147,-0.327) (0.153,-0.245) (0.160,-0.327) (0.167,-0.408) (0.173,-0.490) (0.180,-0.572) (0.187,-0.490) (0.193,-0.572) (0.200,-0.490) (0.207,-0.408) (0.213,-0.327) (0.220,-0.408) (0.227,-0.490) (0.233,-0.408) (0.240,-0.490) (0.247,-0.408) (0.253,-0.327) (0.260,-0.245) (0.267,-0.327) (0.273,-0.245) (0.280,-0.163) (0.287,-0.082) (0.293,-0.163) (0.300,-0.245) (0.307,-0.163) (0.313,-0.245) (0.320,-0.327) (0.327,-0.408) (0.333,-0.327) (0.340,-0.408) (0.347,-0.327) (0.353,-0.408) (0.360,-0.490) (0.367,-0.408) (0.373,-0.327) (0.380,-0.245) (0.387,-0.163) (0.393,-0.245) (0.400,-0.327) (0.407,-0.245) (0.413,-0.163) (0.420,-0.082) (0.427,-0.163) (0.433,-0.245) (0.440,-0.163) (0.447,-0.245) (0.453,-0.327) (0.460,-0.245) (0.467,-0.327) (0.473,-0.245) (0.480,-0.327) (0.487,-0.408) (0.493,-0.327) (0.500,-0.408) (0.507,-0.490) (0.513,-0.572) (0.520,-0.653) (0.527,-0.735) (0.533,-0.653) (0.540,-0.572) (0.547,-0.653) (0.553,-0.735) (0.560,-0.653) (0.567,-0.572) (0.573,-0.653) (0.580,-0.735) (0.587,-0.816) (0.593,-0.898) (0.600,-0.816) (0.607,-0.735) (0.613,-0.653) (0.620,-0.735) (0.627,-0.653) (0.633,-0.572) (0.640,-0.653) (0.647,-0.572) (0.653,-0.653) (0.660,-0.735) (0.667,-0.653) (0.673,-0.572) (0.680,-0.653) (0.687,-0.572) (0.693,-0.490) (0.700,-0.572) (0.707,-0.653) (0.713,-0.572) (0.720,-0.490) (0.727,-0.408) (0.733,-0.490) (0.740,-0.572) (0.747,-0.490) (0.753,-0.572) (0.760,-0.490) (0.767,-0.408) (0.773,-0.490) (0.780,-0.572) (0.787,-0.490) (0.793,-0.572) (0.800,-0.490) (0.807,-0.408) (0.813,-0.327) (0.820,-0.245) (0.827,-0.327) (0.833,-0.245) (0.840,-0.327) (0.847,-0.408) (0.853,-0.327) (0.860,-0.245) (0.867,-0.327) (0.873,-0.408) (0.880,-0.327) (0.887,-0.245) (0.893,-0.163) (0.900,-0.082) (0.907,-0.163) (0.913,-0.245) (0.920,-0.163) (0.927,-0.245) (0.933,-0.327) (0.940,-0.408) (0.947,-0.490) (0.953,-0.572) (0.960,-0.490) (0.967,-0.408) (0.973,-0.490) (0.980,-0.572) (0.987,-0.490) (0.993,-0.572) (1.000,-0.490)};
  \addplot[thick,color=thmcol] coordinates {(0.000,0.000) (0.007,-0.106) (0.013,-0.204) (0.020,-0.309) (0.027,-0.230) (0.033,-0.259) (0.040,-0.339) (0.047,-0.431) (0.053,-0.397) (0.060,-0.483) (0.067,-0.587) (0.073,-0.537) (0.080,-0.635) (0.087,-0.661) (0.093,-0.661) (0.100,-0.698) (0.107,-0.702) (0.113,-0.593) (0.120,-0.635) (0.127,-0.738) (0.133,-0.888) (0.140,-0.905) (0.147,-0.933) (0.153,-0.912) (0.160,-0.950) (0.167,-0.989) (0.173,-1.048) (0.180,-1.090) (0.187,-1.077) (0.193,-1.108) (0.200,-1.100) (0.207,-0.945) (0.213,-0.905) (0.220,-1.034) (0.227,-0.893) (0.233,-0.864) (0.240,-0.941) (0.247,-0.867) (0.253,-0.866) (0.260,-0.916) (0.267,-0.968) (0.273,-1.049) (0.280,-1.045) (0.287,-0.957) (0.293,-0.984) (0.300,-0.950) (0.307,-0.796) (0.313,-0.896) (0.320,-0.875) (0.327,-0.900) (0.333,-0.986) (0.340,-1.070) (0.347,-1.071) (0.353,-1.036) (0.360,-1.079) (0.367,-1.075) (0.373,-1.016) (0.380,-0.991) (0.387,-0.921) (0.393,-0.883) (0.400,-0.843) (0.407,-0.694) (0.413,-0.643) (0.420,-0.632) (0.427,-0.574) (0.433,-0.650) (0.440,-0.677) (0.447,-0.612) (0.453,-0.561) (0.460,-0.434) (0.467,-0.434) (0.473,-0.553) (0.480,-0.394) (0.487,-0.305) (0.493,-0.391) (0.500,-0.279) (0.507,-0.276) (0.513,-0.424) (0.520,-0.459) (0.527,-0.500) (0.533,-0.370) (0.540,-0.435) (0.547,-0.456) (0.553,-0.478) (0.560,-0.509) (0.567,-0.584) (0.573,-0.566) (0.580,-0.478) (0.587,-0.427) (0.593,-0.503) (0.600,-0.597) (0.607,-0.587) (0.613,-0.645) (0.620,-0.696) (0.627,-0.833) (0.633,-0.674) (0.640,-0.599) (0.647,-0.679) (0.653,-0.605) (0.660,-0.495) (0.667,-0.690) (0.673,-0.735) (0.680,-0.767) (0.687,-0.714) (0.693,-0.723) (0.700,-0.742) (0.707,-0.747) (0.713,-0.596) (0.720,-0.419) (0.727,-0.462) (0.733,-0.538) (0.740,-0.318) (0.747,-0.398) (0.753,-0.445) (0.760,-0.442) (0.767,-0.402) (0.773,-0.318) (0.780,-0.286) (0.787,-0.358) (0.793,-0.316) (0.800,-0.296) (0.807,-0.143) (0.813,-0.144) (0.820,-0.253) (0.827,-0.338) (0.833,-0.220) (0.840,-0.264) (0.847,-0.436) (0.853,-0.483) (0.860,-0.483) (0.867,-0.386) (0.873,-0.469) (0.880,-0.415) (0.887,-0.350) (0.893,-0.407) (0.900,-0.422) (0.907,-0.278) (0.913,-0.137) (0.920,-0.067) (0.927,-0.040) (0.933,0.053) (0.940,0.041) (0.947,0.033) (0.953,-0.037) (0.960,-0.036) (0.967,-0.043) (0.973,0.183) (0.980,0.168) (0.987,0.271) (0.993,0.379) (1.000,0.364)};
  \legend{{$\pm1$ steps},{Gaussian steps}}
\end{axis}
\end{tikzpicture}
\caption{Normalised random walks $B_n(t)=S_{\lfloor nt\rfloor}/\sqrt n$ (two independent simulations, seeds fixed). Already at
$n=150$ the Bernoulli and the Gaussian walk are visually indistinguishable, the point of the animation \file{lec14_4}
(there up to $n=2000$) \ts{24}{14}{21:46}\ts{24}{14}{22:47}.}\label{ch11:fig-donsker}
\end{figure}

\begin{intuition}[Universality: the CLT as a fixed point]
Only the first two moments of the steps survive in the limit; all the rest (skewness, kurtosis, lattice structure) is
washed out at rate $1/\sqrt n$. This is the renormalisation-group picture: $S_n/\sqrt n$ is the coarse-graining of
$S_n$ by a factor $n$, and the Gaussian is the fixed point of this flow, with $\sigma^2$ the only relevant coupling. The
consequence for modelling is what the lecture emphasised: one may use Brownian motion when the microscopic steps are
non-Gaussian, provided they have finite variance and are not too dependent \ts{24}{14}{22:47}\ts{24}{14}{23:52}. The 2013
lecture gives the two physical examples: the pollen grain receives a huge number of tiny left/right pushes from
water molecules (Einstein), and a stock price is pushed up or down by a huge number of tiny buy and sell orders
\ts{13}{17}{17:19}\ts{13}{17}{18:21}\ts{13}{17}{21:30}\ts{13}{17}{23:45}. The limit fails when the variance is infinite (stable L\'evy
processes with jumps) or when the steps are strongly dependent (long memory, fractional Brownian motion).
\end{intuition}

\begin{finance}[Sampling frequency and the model]
An analyst observes prices on a grid of spacing $h$: if the increments are like a random walk, the sampled process at fine
resolution looks like Brownian motion (\cref{ch11:thm-donsker} with $n\sim1/h$) \ts{13}{17}{17:19}. Two warnings. At very short horizons
(ticks) market microstructure (bid--ask bounce) makes the increments negatively correlated, so the limit is not
Brownian; and price jumps have heavy tails, so ``finite variance'' can fail (\cref{ch03:sec-heavy}). Universality is a
statement about the \emph{long-horizon} behaviour of the sum of many small steps.
\end{finance}

\paragraph{Where the theorem is used.} Because $F(B_n)\to F(B)$ for continuous $F$, asymptotic distributions of statistics
of random walks are the distributions of functionals of Brownian motion. For the simple random walk one even has an exact
discrete analogue of the maximum formula of \cref{ch11:sec-reflection}: $\Prob[\max_{k\le n}S_k\ge a]=2\Prob[S_n>a]+\Prob[S_n=a]$
(checked exactly for $n=100$, $a=10$: both sides equal $0.3197$), so $\Prob[\max_{k\le n}S_k\ge a\sqrt n]\to2(1-\Phi(a))$, the
Brownian result. Similarly $n^{-2}\sum_kS_{k-1}^2\to\int_0^1B(t)^2\dd t$ and $n^{-1}\sum_kS_{k-1}X_k\to\int_0^1B\dd B$,
which is how the Dickey--Fuller limit of \cref{ch09:eq-df} (\cref{ch09:sec-df}) is obtained; the Kolmogorov--Smirnov statistic
of \cref{ch11:sec-bridge} is another example.

% ==================================================================
\section{Sample paths: roughness and quadratic variation}\label{ch11:sec-paths}

\subsection{Size, zeros and the envelope}
Almost every Brownian path has the following features (2013 lecture): it \emph{crosses the $t$-axis infinitely often}, in
fact it has zeros accumulating at $t=0$ and at every later zero, and it neither drifts off nor stays near zero; at time $t_0$
it is $N(0,t_0)$, so it stays ``around'' the curves $\pm\sqrt t$ (typical size $\sqrt{t_0}$, never $100$ times that in
practice) \ts{13}{17}{25:48}\ts{13}{17}{26:48}. The precise version of ``it does not deviate too much'' is the law of the
iterated logarithm (not proved in the course): $\limsup_{t\to\infty}B(t)/\sqrt{2t\ln\ln t}=1$ almost surely, so the
$\sqrt t$ envelope is exceeded infinitely often, but only by slowly growing factors (compare \cref{ch04:sec-rw}).

\subsection{Non-differentiability}
The path is continuous by definition, but it is nowhere differentiable, which was ``really surprising'' in the words of the 2013 lecture,
and matters because the whole of classical calculus is built on derivatives \ts{13}{17}{27:48}\ts{13}{17}{28:53}.

\begin{proposition}[No derivative]\label{ch11:prop-nodiff}
Fix $t_0\ge0$. With probability one, $B$ is not differentiable at $t_0$. In fact, with probability one there is no
$t_0$ at which $B$ is differentiable (Paley--Wiener--Zygmund, 1933).
\end{proposition}
\begin{proof}[Proof of the first statement]
The heuristic of the slides \ts{24}{14}{1:02:33}\ts{24}{14}{1:03:39}: the difference quotient $D(\Delta t)=[B(t_0+\Delta t)-B(t_0)]/\Delta t$
is $N(0,1/\Delta t)$, so for every fixed $M<\infty$, $\Prob(|D|\le M)=\Prob(|N(0,1)|\le M\sqrt{\Delta t})\le M\sqrt{2\Delta t/\pi}\to0$:
the slope does not settle down; it is of order $(\Delta t)^{-1/2}$ (equivalently $\Delta B=O_P(\sqrt{\Delta t})\gg\Delta t$).
To make this a proof, note that a finite derivative at $t_0$ forces, for some integers $K,m$, $|B(t_0+s)-B(t_0)|\le Ks$ for all
$0<s<1/m$; the event $E_{K,m}$ so defined satisfies $\Prob(E_{K,m})\le\Prob(|B(t_0+\delta)-B(t_0)|\le K\delta)\le K\sqrt{2\delta/\pi}$
for every $\delta<1/m$, hence $\Prob(E_{K,m})=0$; take the countable union over $K,m$ \ts{13}{17}{43:15}\ts{13}{17}{44:19}\ts{13}{17}{45:23}.
The statement for all $t_0$ simultaneously needs a covering argument and is quoted.
\end{proof}

\begin{coursenote}[Erratum in the 2013 proof of non-differentiability]
The written proof (\file{lecnote17}, Proposition~2.4) bounds the event by $\Prob(M(\delta)\le A\delta)$ for the maximum
$M(\delta)=\max_{s\le\delta}B(s)$, and computes it as ``$2(1-\Phi(A\sqrt\delta))$, which tends to zero''. That value tends to $1$, not
$0$. By \cref{ch11:prop-max} the correct value is $\Prob(M(\delta)\le A\delta)=2\Phi(A\sqrt\delta)-1$, which does tend to $0$; the
lecture itself stumbles over the direction of the inequality at this point \ts{13}{17}{46:23}. Both versions are statements for a fixed $t_0$; the slide of 2024 (``not differentiable at $t$, for any $t\ge0$'') is likewise a statement for each fixed $t$.
The argument above avoids the maximum altogether.
\end{coursenote}

\subsection{Quadratic variation}
Let $\Pi=\{0=t_0<t_1<\dots<t_n=T\}$ be a partition of $[0,T]$ with mesh $|\Pi|=\max_j(t_j-t_{j-1})$ and define the
$\Pi$-quadratic variation
\ts{24}{14}{1:03:39}\ts{24}{14}{1:04:44}\ts{24}{14}{1:05:45}
\[
  QV_\Pi([0,T])=\sum_{j=1}^{n}\bigl[B(t_j)-B(t_{j-1})\bigr]^2 .
\]

\begin{theorem}[Quadratic variation of Brownian motion]\label{ch11:thm-qv}
$QV_\Pi([0,T])\to T$ as $|\Pi|\to0$ in $L^2$ (hence in probability), and almost surely along any sequence of partitions with
$\sum_k|\Pi_k|<\infty$, for instance the dyadic partitions \ts{24}{14}{1:06:57}\ts{13}{17}{47:24}.
\end{theorem}
\begin{proof}
Write $\Delta_j=t_j-t_{j-1}$ and $\Delta_jB=B(t_j)-B(t_{j-1})\sim N(0,\Delta_j)$, independent over $j$. Then
$\E[(\Delta_jB)^2]=\Delta_j$, $\Var[(\Delta_jB)^2]=3\Delta_j^2-\Delta_j^2=2\Delta_j^2$ (the fourth moment of $N(0,\Delta_j)$ is $3\Delta_j^2$), so
\[
  \E[Q]V_\Pi=\sum_j\Delta_j=T,\qquad
  \Var\,QV_\Pi=2\sum_j\Delta_j^2\le2|\Pi|\sum_j\Delta_j=2|\Pi|\,T\xrightarrow[|\Pi|\to0]{}0 .
\]
For equally spaced points $QV_\Pi\eqd(T/n)\,\chi^2_n$ \ts{24}{14}{1:08:03}, mean $T$ and variance $2T^2/n$. Chebyshev's inequality gives
convergence in probability; if $\sum_k|\Pi_k|<\infty$ then $\sum_k\Prob(|QV_{\Pi_k}-T|>\varepsilon)\le\sum_k2|\Pi_k|T/\varepsilon^2<\infty$ and
Borel--Cantelli gives almost sure convergence.
\end{proof}

\begin{coursenote}[Erratum and remarks on the proof in the notes and slides]
\begin{itemize}
  \item The 2013 proof invokes the law of large numbers for the sum of $n$ i.i.d.\ variables $(B(t_{i+1})-B(t_i))^2$ of mean $T/n$; the
      summands change with $n$ (a triangular array), so the classical law does not apply verbatim, and the mean and variance computation above (Chebyshev)
      is the clean way; the quantity converges in $L^2$ and, for nested partitions, almost surely. The lecture's own computation
      had a slip corrected by a student (the mean of $X_i^2$ is $T/n$, not $T^2/n^2$) \ts{13}{17}{56:07}.
  \item Slide~19 of \file{lec14\_1} (PDF pages 73--74) says ``$QV([0,T])=1$ for virtually every path''; the value is $T$ (it equals $1$ only for $T=1$), as the
      claim on the same slide states.
  \item On the slide of the Brownian bridge (slide~13, \cref{ch11:sec-bridge}) ``Moments of $B(t)$'' should read
      ``moments of $X(t)$''.
  \item Slide~2 dates Einstein's diffusion paper to 1904; it appeared in 1905 (as in the 2013 note and in
      our Motivation).
  \item The diffusion equation on slide~6 (PDF page 15) is written $\partial_tp=\tfrac12\partial_y^2p$, which is the case $\sigma=1$ announced there
      (``easy to verify for $s=0$ and $\sigma=1$''); for a general diffusion coefficient the factor is $\tfrac12\sigma^2$, \cref{ch11:prop-heat}.
\end{itemize}
\end{coursenote}

The theorem says that the \emph{sum of squared increments is not random and does not depend on the path}: $[B]_T=T$ for
almost every path \ts{24}{14}{1:06:57}\ts{24}{14}{1:09:17}. Compare a continuously differentiable function $f$: by the mean value theorem
$\sum_j(f(t_j)-f(t_{j-1}))^2=\sum_jf'(s_j)^2\Delta_j^2\le\max|f'|^2\,|\Pi|\,T\to0$, so its quadratic variation is zero
\ts{13}{17}{48:29}\ts{13}{17}{49:39}\ts{13}{17}{50:42}\ts{13}{17}{52:44}. Brownian paths ``bounce back and forth too much'': the squared moves
do not vanish as the grid is refined. Two corollaries.

\begin{corollary}[Infinite total variation]\label{ch11:cor-tv}
Almost surely $\sum_j|B(t_j)-B(t_{j-1})|\to\infty$ as $|\Pi|\to0$.
\end{corollary}
\begin{proof}
$\sum_j\Delta_jB^2\le\max_j|\Delta_jB|\cdot\sum_j|\Delta_jB|$. The left side tends to $T>0$ and, by uniform continuity of the
path, $\max_j|\Delta_jB|\to0$; hence the sum of absolute increments diverges.
\end{proof}
So one cannot define $\int f\dd B$ pathwise as a Riemann--Stieltjes integral, and the stochastic (It\^o) integral of \cref{ch:stochcalc} is needed.

The rule of thumb is written
\begin{equation}\label{ch11:eq-dB2}
  (\dd B)^2=\dd t,\qquad \dd B\,\dd t=0,\qquad (\dd t)^2=0,
\end{equation}
where $\dd B\sim\sqrt{\dd t}$ carries the whole effect at second order: the increments of a Brownian path are of order
$\sqrt{\Delta t}$, their squares of order $\Delta t$ (summing to a finite number), while a product $\Delta B\,\Delta t$
is $O(\Delta t^{3/2})$ and vanishes on summation \ts{13}{17}{53:49}. This is the reason why the Taylor expansion of a function of $B$ has a second-order term (\cref{ch11:sec-gbm}).

\begin{figure}[ht]
\centering
\begin{tikzpicture}
\begin{axis}[width=0.95\linewidth,height=5cm,xlabel={$t$},ylabel={$QV_n(t)$},xmin=0,xmax=1,ymin=0,ymax=1.4,
  legend style={at={(0.5,1.03)},anchor=south,legend columns=3,font=\footnotesize},grid=major,grid style={black!10}]
  \addplot[thick,color=defcol] coordinates {(0.000,0.000) (0.000,0.000) (0.020,0.083) (0.020,0.083) (0.040,0.214) (0.040,0.214) (0.060,0.217) (0.060,0.217) (0.080,0.224) (0.080,0.224) (0.100,0.228) (0.100,0.228) (0.120,0.229) (0.120,0.229) (0.140,0.311) (0.140,0.311) (0.160,0.312) (0.160,0.312) (0.180,0.327) (0.180,0.327) (0.200,0.547) (0.200,0.547) (0.220,0.548) (0.220,0.548) (0.240,0.551) (0.240,0.551) (0.260,0.552) (0.260,0.552) (0.280,0.561) (0.280,0.561) (0.300,0.584) (0.300,0.584) (0.320,0.587) (0.320,0.587) (0.340,0.591) (0.340,0.591) (0.360,0.593) (0.360,0.593) (0.380,0.611) (0.380,0.611) (0.400,0.612) (0.400,0.612) (0.420,0.612) (0.420,0.612) (0.440,0.659) (0.440,0.659) (0.460,0.665) (0.460,0.665) (0.480,0.671) (0.480,0.671) (0.500,0.671) (0.500,0.671) (0.520,0.677) (0.520,0.677) (0.540,0.752) (0.540,0.752) (0.560,0.753) (0.560,0.753) (0.580,0.755) (0.580,0.755) (0.600,0.775) (0.600,0.775) (0.620,0.790) (0.620,0.790) (0.640,0.792) (0.640,0.792) (0.660,0.808) (0.660,0.808) (0.680,0.814) (0.680,0.814) (0.700,0.815) (0.700,0.815) (0.720,0.824) (0.720,0.824) (0.740,0.984) (0.740,0.984) (0.760,1.004) (0.760,1.004) (0.780,1.023) (0.780,1.023) (0.800,1.078) (0.800,1.078) (0.820,1.080) (0.820,1.080) (0.840,1.090) (0.840,1.090) (0.860,1.094) (0.860,1.094) (0.880,1.117) (0.880,1.117) (0.900,1.117) (0.900,1.117) (0.920,1.117) (0.920,1.117) (0.940,1.157) (0.940,1.157) (0.960,1.168) (0.960,1.168) (0.980,1.168) (0.980,1.168) (1.000,1.193)};
  \addplot[thick,color=thmcol] coordinates {(0.000,0.000) (0.010,0.009) (0.020,0.025) (0.030,0.038) (0.040,0.045) (0.050,0.052) (0.060,0.058) (0.070,0.065) (0.080,0.083) (0.090,0.097) (0.100,0.103) (0.110,0.110) (0.120,0.116) (0.130,0.133) (0.140,0.146) (0.150,0.151) (0.160,0.154) (0.170,0.171) (0.180,0.180) (0.190,0.193) (0.200,0.199) (0.210,0.216) (0.220,0.224) (0.230,0.234) (0.240,0.245) (0.250,0.249) (0.260,0.260) (0.270,0.272) (0.280,0.280) (0.290,0.288) (0.300,0.307) (0.310,0.313) (0.320,0.323) (0.330,0.328) (0.340,0.335) (0.350,0.344) (0.360,0.352) (0.370,0.362) (0.380,0.385) (0.390,0.402) (0.400,0.410) (0.410,0.423) (0.420,0.433) (0.430,0.438) (0.440,0.452) (0.450,0.464) (0.460,0.474) (0.470,0.483) (0.480,0.494) (0.490,0.505) (0.500,0.509) (0.510,0.515) (0.520,0.526) (0.530,0.530) (0.540,0.537) (0.550,0.552) (0.560,0.558) (0.570,0.574) (0.580,0.583) (0.590,0.598) (0.600,0.607) (0.610,0.623) (0.620,0.634) (0.630,0.645) (0.640,0.657) (0.650,0.664) (0.660,0.674) (0.670,0.687) (0.680,0.695) (0.690,0.702) (0.700,0.713) (0.710,0.729) (0.720,0.738) (0.730,0.745) (0.740,0.753) (0.750,0.764) (0.760,0.776) (0.770,0.789) (0.780,0.800) (0.790,0.812) (0.800,0.822) (0.810,0.831) (0.820,0.841) (0.830,0.850) (0.840,0.853) (0.850,0.874) (0.860,0.881) (0.870,0.882) (0.880,0.896) (0.890,0.898) (0.900,0.911) (0.910,0.925) (0.920,0.933) (0.930,0.939) (0.940,0.953) (0.950,0.963) (0.960,0.968) (0.970,0.988) (0.980,0.997) (0.990,1.013) (1.000,1.031)};
  \addplot[dashed,black!70] coordinates {(0,0) (1,1)};
  \legend{{$n=50$ steps},{$n=1000$ steps},{$T=t$}}
\end{axis}
\end{tikzpicture}
\caption{Cumulative quadratic variation $QV_n(t)=\sum_{k\le nt}(\Delta B_k)^2$ of a Gaussian random walk
$\Delta B_k=X_k/\sqrt n$, at two resolutions. At $t=1$ the values are $1.19$ and $1.03$ (theory: standard deviation
$\sqrt{2/n}$, i.e.\ $0.20$ and $0.045$); the path converges to the straight line $t$, as in \file{lec14\_9}
(there $n$ up to $2000$) \ts{24}{14}{1:08:03}\ts{24}{14}{1:09:17}.}\label{ch11:fig-qv}
\end{figure}

\begin{finance}[Realised variance and stochastic volatility]
If the log price is $\log(P_t/P_0)=\mu t+\sigma B(t)$, the sum of squared log returns over a period $[0,T]$ (the \emph{realised
variance}) estimates $\sigma^2T$ \emph{without knowing the drift and without error as the sampling becomes finer}: with
$n$ returns its relative standard deviation is $\sqrt{2/n}$ (about $8.9\%$ for $n=252$ daily returns of one year, so
$8.9\%/2\approx4.5\%$ in volatility terms). This is the theoretical basis for measuring volatility by high-frequency data
(\cref{ch:volatility}). The lecture proposes plotting the cumulative squared returns of an asset against time: for a Brownian model the
graph is a straight line of slope $\sigma^2$; slope changes and curvature reveal periods of low and high volatility and
suggest a model of dynamic volatility \ts{24}{14}{1:10:21}\ts{24}{14}{1:11:22}\ts{24}{14}{1:12:22}. Drift, by contrast, cannot be
learned by refining the sampling: $\hat\mu$ has standard error $\sigma/\sqrt T$ whatever $n$ is (2024 PS4 Problem~3,
\cref{ch:volatility}).
\end{finance}


% ==================================================================
\section{The reflection principle, the maximum and hitting times}\label{ch11:sec-reflection}

Suppose Brownian motion models a stock price during a trading day (9:30 to 16:00). Traders care about the \emph{highest
and lowest} values reached, not only the closing value \ts{13}{17}{31:00}: barrier options pay if a level is touched, a stop-loss triggers when it is
reached, a firm defaults when its value hits a threshold. The two objects are the running maximum and the first hitting time.
For a standard Brownian motion $B$ and levels $a>0$ define
\begin{equation}\label{ch11:eq-defM}
  M(t)=\max_{0\le u\le t}B(u),\qquad \tau_a=\min\{u\ge0:\ B(u)=a\}.
\end{equation}
The maximum exists because a continuous function attains its maximum on the compact set $[0,t]$ \ts{24}{14}{28:05}\ts{13}{17}{31:00}.
Since the path is continuous and starts at $0<a$, hitting $a$ before $t$ and having a maximum $\ge a$ are the same event:
\begin{equation}\label{ch11:eq-tauM}
  \{\tau_a\le t\}=\{M(t)\ge a\}.
\end{equation}
$\tau_a$ is a \emph{stopping time} (\cref{ch04:def-stop}): whether $\tau_a\le t$ is decided by the path up to time $t$
\ts{24}{14}{24:53}\ts{13}{17}{33:04}.

\begin{lemma}[Reflection principle]\label{ch11:lem-reflection}
Let $\widetilde B(u)=B(u)$ for $u\le\tau_a$ and $\widetilde B(u)=a-\bigl(B(u)-a\bigr)=2a-B(u)$ for $u>\tau_a$: the path is
reflected in the level $a$ from the first hitting time on. Then $\widetilde B$ is again a standard Brownian motion
\ts{24}{14}{24:53}\ts{24}{14}{25:57}.
\end{lemma}
The proof uses the \emph{strong Markov property}: at the stopping time $\tau_a$ the process $B(\tau_a+u)-B(\tau_a)$ is a new
Brownian motion independent of the past $\mathcal F_{\tau_a}$; the Markov property of \cref{ch11:prop-elem}(ii) still holds when the fixed time $s$ is replaced by a
stopping time, because $\tau_a$ can be approximated by discrete stopping times (\cref{ch04:thm-stopped}) (2013 note: ``the
distribution of $B(t)-B(\tau_a)$ is not affected by the fact that we conditioned on $\tau_a<t$'') \ts{13}{17}{35:54}. Since $-B$ is
also a Brownian motion, gluing the old path up to $\tau_a$ to the reflected new increment $-(B(\tau_a+u)-a)$ gives
a process with the same law. The lemma is the statement that after touching the level the path is equally likely to go
up or down \ts{24}{14}{24:53}\ts{24}{14}{25:57}\ts{24}{14}{26:57}; \cref{ch11:fig-refl} shows one path and its reflection (the deck
\file{lec14\_5} shows it for $x=50$ and $10\,000$ steps).

\begin{figure}[ht]
\centering
\begin{tikzpicture}
\begin{axis}[width=0.95\linewidth,height=5.4cm,xlabel={$t$},xmin=0,xmax=1,ymin=-0.7,ymax=1.5,
  legend style={at={(0.5,1.03)},anchor=south,legend columns=3,font=\footnotesize},grid=major,grid style={black!10}]
  \addplot[thick,color=defcol] coordinates {(0.000,0.000) (0.005,-0.042) (0.010,0.003) (0.015,0.076) (0.020,0.149) (0.025,0.277) (0.030,0.250) (0.035,0.289) (0.040,0.263) (0.045,0.162) (0.050,0.112) (0.055,0.122) (0.060,0.057) (0.065,0.044) (0.070,0.123) (0.075,0.163) (0.080,0.204) (0.085,0.228) (0.090,0.216) (0.095,0.084) (0.100,0.154) (0.105,0.047) (0.110,0.062) (0.115,0.055) (0.120,0.065) (0.125,0.082) (0.130,0.059) (0.135,0.122) (0.140,0.032) (0.145,0.088) (0.150,-0.032) (0.155,0.052) (0.160,0.016) (0.165,0.042) (0.170,0.149) (0.175,-0.004) (0.180,-0.026) (0.185,0.014) (0.190,0.078) (0.195,0.174) (0.200,0.217) (0.205,0.253) (0.210,0.261) (0.215,0.279) (0.220,0.284) (0.225,0.355) (0.230,0.366) (0.235,0.413) (0.240,0.527) (0.245,0.610) (0.250,0.586) (0.255,0.579) (0.260,0.631) (0.265,0.585) (0.270,0.456) (0.275,0.384) (0.280,0.422) (0.285,0.407) (0.290,0.382) (0.295,0.435) (0.300,0.361) (0.305,0.521) (0.310,0.522) (0.315,0.499) (0.320,0.643) (0.325,0.566) (0.330,0.531) (0.335,0.525) (0.340,0.457) (0.345,0.433) (0.350,0.474) (0.355,0.486) (0.360,0.483) (0.365,0.553) (0.370,0.686) (0.375,0.797) (0.380,0.648) (0.385,0.745) (0.390,0.800) (0.395,0.829) (0.400,0.760) (0.405,0.728) (0.410,0.713) (0.415,0.614) (0.420,0.702) (0.425,0.765) (0.430,0.764) (0.435,0.697) (0.440,0.792) (0.445,0.837) (0.450,0.912) (0.455,0.883) (0.460,1.052) (0.465,1.042) (0.470,1.181) (0.475,1.039) (0.480,1.061) (0.485,1.066) (0.490,0.888) (0.495,0.842) (0.500,0.756) (0.505,0.656) (0.510,0.603) (0.515,0.520) (0.520,0.543) (0.525,0.576) (0.530,0.595) (0.535,0.537) (0.540,0.502) (0.545,0.621) (0.550,0.597) (0.555,0.595) (0.560,0.554) (0.565,0.591) (0.570,0.654) (0.575,0.710) (0.580,0.740) (0.585,0.790) (0.590,0.775) (0.595,0.819) (0.600,0.749) (0.605,0.756) (0.610,0.817) (0.615,0.746) (0.620,0.801) (0.625,0.823) (0.630,0.737) (0.635,0.616) (0.640,0.499) (0.645,0.425) (0.650,0.448) (0.655,0.465) (0.660,0.552) (0.665,0.713) (0.670,0.679) (0.675,0.535) (0.680,0.429) (0.685,0.286) (0.690,0.237) (0.695,0.199) (0.700,0.289) (0.705,0.182) (0.710,0.164) (0.715,0.146) (0.720,0.137) (0.725,0.096) (0.730,0.245) (0.735,0.232) (0.740,0.235) (0.745,0.284) (0.750,0.300) (0.755,0.425) (0.760,0.339) (0.765,0.325) (0.770,0.289) (0.775,0.283) (0.780,0.346) (0.785,0.303) (0.790,0.164) (0.795,0.168) (0.800,0.094) (0.805,0.137) (0.810,0.153) (0.815,0.301) (0.820,0.339) (0.825,0.280) (0.830,0.305) (0.835,0.381) (0.840,0.386) (0.845,0.384) (0.850,0.365) (0.855,0.336) (0.860,0.270) (0.865,0.356) (0.870,0.384) (0.875,0.282) (0.880,0.358) (0.885,0.360) (0.890,0.376) (0.895,0.387) (0.900,0.471) (0.905,0.314) (0.910,0.312) (0.915,0.352) (0.920,0.343) (0.925,0.226) (0.930,0.137) (0.935,0.152) (0.940,0.253) (0.945,0.116) (0.950,0.075) (0.955,0.098) (0.960,0.116) (0.965,0.156) (0.970,0.239) (0.975,0.295) (0.980,0.193) (0.985,0.174) (0.990,0.098) (0.995,0.100) (1.000,0.209)};
  \addplot[thick,dashed,color=thmcol] coordinates {(0.390,0.800) (0.390,0.800) (0.395,0.771) (0.400,0.840) (0.405,0.872) (0.410,0.887) (0.415,0.986) (0.420,0.898) (0.425,0.835) (0.430,0.836) (0.435,0.903) (0.440,0.808) (0.445,0.763) (0.450,0.688) (0.455,0.717) (0.460,0.548) (0.465,0.558) (0.470,0.419) (0.475,0.561) (0.480,0.539) (0.485,0.534) (0.490,0.712) (0.495,0.758) (0.500,0.844) (0.505,0.944) (0.510,0.997) (0.515,1.080) (0.520,1.057) (0.525,1.024) (0.530,1.005) (0.535,1.063) (0.540,1.098) (0.545,0.979) (0.550,1.003) (0.555,1.005) (0.560,1.046) (0.565,1.009) (0.570,0.946) (0.575,0.890) (0.580,0.860) (0.585,0.810) (0.590,0.825) (0.595,0.781) (0.600,0.851) (0.605,0.844) (0.610,0.783) (0.615,0.854) (0.620,0.799) (0.625,0.777) (0.630,0.863) (0.635,0.984) (0.640,1.101) (0.645,1.175) (0.650,1.152) (0.655,1.135) (0.660,1.048) (0.665,0.887) (0.670,0.921) (0.675,1.065) (0.680,1.171) (0.685,1.314) (0.690,1.363) (0.695,1.401) (0.700,1.311) (0.705,1.418) (0.710,1.436) (0.715,1.454) (0.720,1.463) (0.725,1.504) (0.730,1.355) (0.735,1.368) (0.740,1.365) (0.745,1.316) (0.750,1.300) (0.755,1.175) (0.760,1.261) (0.765,1.275) (0.770,1.311) (0.775,1.317) (0.780,1.254) (0.785,1.297) (0.790,1.436) (0.795,1.432) (0.800,1.506) (0.805,1.463) (0.810,1.447) (0.815,1.299) (0.820,1.261) (0.825,1.320) (0.830,1.295) (0.835,1.219) (0.840,1.214) (0.845,1.216) (0.850,1.235) (0.855,1.264) (0.860,1.330) (0.865,1.244) (0.870,1.216) (0.875,1.318) (0.880,1.242) (0.885,1.240) (0.890,1.224) (0.895,1.213) (0.900,1.129) (0.905,1.286) (0.910,1.288) (0.915,1.248) (0.920,1.257) (0.925,1.374) (0.930,1.463) (0.935,1.448) (0.940,1.347) (0.945,1.484) (0.950,1.525) (0.955,1.502) (0.960,1.484) (0.965,1.444) (0.970,1.361) (0.975,1.305) (0.980,1.407) (0.985,1.426) (0.990,1.502) (0.995,1.500) (1.000,1.391)};
  \addplot[thin,black!60] coordinates {(0,0.8) (1,0.8)};
  \addplot[thin,dotted,black!60] coordinates {(0.390,-0.7) (0.390,0.8)};
  \node[font=\footnotesize,anchor=north west] at (axis cs:0.395,-0.42) {$\tau_a$};
  \node[font=\footnotesize,anchor=south west] at (axis cs:0.01,0.82) {level $a=0.8$};
  \legend{{path $B$},{reflected path $\widetilde B$ after $\tau_a$}}
\end{axis}
\end{tikzpicture}
\caption{A simulated Brownian path ($200$ steps, $a=0.8$) and its reflection in the level $a$ after the first hitting time
$\tau_a$. The two paths have the same probability; the original ends below $a$, its reflection above.}\label{ch11:fig-refl}
\end{figure}

\begin{proposition}[Distribution of the maximum]\label{ch11:prop-max}
For $t>0$ and $a>0$,
\begin{equation}\label{ch11:eq-max}
  \Prob[M(t)\ge a]=2\,\Prob[B(t)>a]=2\Bigl[1-\Phi\Bigl(\frac a{\sqrt t}\Bigr)\Bigr].
\end{equation}
Equivalently $M(t)\eqd|B(t)|$, with density $\sqrt{2/(\pi t)}\,e^{-m^2/(2t)}$ on $m>0$, and $\E[M(t)]=\sqrt{2t/\pi}$.
\end{proposition}
\begin{proof}
Two views, both used in the lectures.
\emph{Path pairing} \ts{24}{14}{28:05}\ts{24}{14}{29:20}\ts{24}{14}{30:27}. Let $A_{\mathrm{END}}=\{B(t)>a\}$ and $A_{\mathrm{MAX}}=\{M(t)\ge a\}$.
By continuity $A_{\mathrm{END}}\subset A_{\mathrm{MAX}}$ (a path ending above $a$ has crossed it, so $\tau_a<t$). Every
$\omega_0\in A_{\mathrm{END}}$ is paired with its reflection $\omega_1$ (\cref{ch11:lem-reflection}), which
ends below $a$ but has still maximum $\ge a$, and has the same probability; conversely every path in $A_{\mathrm{MAX}}\setminus A_{\mathrm{END}}$
arises this way. Hence $\Prob[A_{\mathrm{MAX}}]=\Prob[A_{\mathrm{MAX}}\cap A_{\mathrm{END}}]+\Prob[A_{\mathrm{MAX}}\cap A_{\mathrm{END}}^c]
=\Prob[A_{\mathrm{END}}]+\Prob[A_{\mathrm{END}}]$ (the event $B(t)=a$ has probability $0$).
\emph{Conditioning on $\tau_a$} \ts{13}{17}{34:09}\ts{13}{17}{36:56}\ts{13}{17}{40:08}. Given $\tau_a<t$, the strong Markov property says $B(t)-B(\tau_a)=B(t)-a$
is symmetric, so $\Prob[B(t)>a\mid\tau_a<t]=\Prob[B(t)<a\mid\tau_a<t]=\tfrac12$. Therefore
$\Prob[M(t)\ge a]=\Prob[\tau_a<t]=2\Prob[B(t)>a,\ \tau_a<t]=2\Prob[B(t)>a]$, the last step because $B(t)>a$ already
implies $\tau_a<t$.
The rest follows from $\Prob[B(t)>a]=1-\Phi(a/\sqrt t)$, by differentiating in $a$, and by $\E[M]=\int_0^\infty\Prob[M>a]\dd a$.
\end{proof}

\begin{remark}[Joint law of maximum and endpoint; not discussed in class]\label{ch11:rem-joint}
The same reflection gives, for $m>0$ and $x\le m$, $\Prob[M(t)\ge m,\ B(t)\le x]=\Prob[B(t)\ge2m-x]=1-\Phi\bigl((2m-x)/\sqrt t\bigr)$
(reflect the part of the path after $\tau_m$: the event $\{B(t)\le x\}$ becomes $\{\widetilde B(t)\ge2m-x\}$).
For $x=m$ this is \eqref{ch11:eq-max}. By symmetry, the minimum satisfies $\Prob[\min_{u\le t}B(u)\le-a]=2[1-\Phi(a/\sqrt t)]$
as well.
\end{remark}

\begin{finance}[Barrier events: a path is twice as likely to touch a level as to end beyond it]
Let $\log(S_t/S_0)=\sigma B(t)$ (no drift, $\sigma=30\%$ per year). The probability that the price is below $S_0e^{-0.2}\approx0.82\,S_0$ \emph{at the end of
one year} is $\Phi(-0.2/0.3)=25.2\%$; the probability that the stock \emph{touches} that level at some time during the year is
$2\Phi(-0.2/0.3)=50.5\%$ (\cref{ch11:prop-max} applied to $-B$). For a knock-out or knock-in
barrier option, or a stop-loss, the touching probability is what matters, and ``double the terminal probability'' is
the first-order rule; a drift or discrete (say daily) monitoring changes the number, the latter by lowering it
(a path can cross between two observations without being seen) \ts{13}{17}{31:00}. Range-based volatility estimators
combine the maximum and the minimum of the log price over a day: with the range $R=\max-\min$ of a driftless Brownian path one has
$\E[R^2]=4\ln2\cdot\sigma^2t$ (Parkinson 1980, derived from the joint law of \cref{ch11:rem-joint}; it is the subject of 2024 PS5 Problem~3 and
of \cref{ch:volatility}); in the 2013 lecture Kempthorne, from the audience, remarks that he uses the range of the process as a more precise measure of volatility than close-to-close prices
\ts{13}{17}{39:05}\ts{13}{17}{40:08}.
\end{finance}

\subsection{First hitting time}
By \eqref{ch11:eq-tauM} and \cref{ch11:prop-max},
\begin{equation}\label{ch11:eq-hitcdf}
  \Prob[\tau_a\le t]=2\Bigl[1-\Phi\Bigl(\frac a{\sqrt t}\Bigr)\Bigr],\qquad
  f_{\tau_a}(t)=\frac{a}{\sqrt{2\pi}}\,t^{-3/2}\,e^{-a^2/(2t)},\quad0<t<\infty,
\end{equation}
\sloppy the density being the derivative in $t$ (using $\frac{\dd}{\dd t}\Phi(a t^{-1/2})=-\tfrac12 a t^{-3/2}\varphi(a t^{-1/2})$)
\ts{24}{14}{31:35}\ts{24}{14}{32:37}\ts{24}{14}{33:44}. This is the L\'evy distribution.

\begin{proposition}[Properties of $\tau_a$]\label{ch11:prop-tau}
\leavevmode
\begin{enumerate}[(i)]
  \item $\tau_a<\infty$ almost surely: a Brownian path reaches every level.
  \item $\E[\tau_a]=\infty$: the tail is $\Prob[\tau_a>t]=2\Phi(a/\sqrt t)-1\sim a\sqrt{2/(\pi t)}$ and $\int^\infty t^{-1/2}\dd t=\infty$.
  \item The density is unimodal with mode at $t=a^2/3$, and its median is $a^2/\bigl(\Phi^{-1}(3/4)\bigr)^2\approx2.198\,a^2$.
  \item Scaling: $\tau_a\eqd a^2\tau_1$.
\end{enumerate}
\end{proposition}
\begin{proof}
(i) $\Prob[\tau_a\le t]\to2[1-\Phi(0)]=1$ as $t\to\infty$.

(ii) as stated.

(iii) $\frac{\dd}{\dd t}\ln f=-\tfrac{3}{2t}+\tfrac{a^2}{2t^2}=0$ at $t=a^2/3$;
the median solves $2[1-\Phi(a/\sqrt t)]=\tfrac12$, i.e.\ $a/\sqrt t=\Phi^{-1}(3/4)=0.6745$.

(iv) from \eqref{ch11:eq-hitcdf}, which depends on $a/\sqrt t$ only.
\end{proof}

\begin{figure}[ht]
\centering
\begin{tikzpicture}
\begin{axis}[width=0.47\linewidth,height=5cm,xlabel={$t$},ylabel={$\Prob[M(t)>x]$},xmin=0,xmax=5,ymin=0,ymax=0.7,
  title={Maximum exceeds $x$ before $t$},title style={font=\footnotesize},
  legend style={at={(0.03,0.97)},anchor=north west,font=\scriptsize},grid=major,grid style={black!10}]
  \addplot[thick,color=defcol] coordinates {(0.040,0.000) (0.141,0.008) (0.242,0.042) (0.344,0.088) (0.445,0.134) (0.546,0.176) (0.647,0.214) (0.749,0.248) (0.850,0.278) (0.951,0.305) (1.052,0.330) (1.153,0.352) (1.255,0.372) (1.356,0.390) (1.457,0.407) (1.558,0.423) (1.660,0.438) (1.761,0.451) (1.862,0.464) (1.963,0.475) (2.064,0.486) (2.166,0.497) (2.267,0.507) (2.368,0.516) (2.469,0.525) (2.571,0.533) (2.672,0.541) (2.773,0.548) (2.874,0.555) (2.976,0.562) (3.077,0.569) (3.178,0.575) (3.279,0.581) (3.380,0.587) (3.482,0.592) (3.583,0.597) (3.684,0.602) (3.785,0.607) (3.887,0.612) (3.988,0.617) (4.089,0.621) (4.190,0.625) (4.291,0.629) (4.393,0.633) (4.494,0.637) (4.595,0.641) (4.696,0.644) (4.798,0.648) (4.899,0.651) (5.000,0.655)};
  \addplot[thick,color=intcol] coordinates {(0.040,0.000) (0.141,0.000) (0.242,0.000) (0.344,0.001) (0.445,0.003) (0.546,0.007) (0.647,0.013) (0.749,0.021) (0.850,0.030) (0.951,0.040) (1.052,0.051) (1.153,0.063) (1.255,0.074) (1.356,0.086) (1.457,0.098) (1.558,0.109) (1.660,0.121) (1.761,0.132) (1.862,0.143) (1.963,0.153) (2.064,0.164) (2.166,0.174) (2.267,0.184) (2.368,0.194) (2.469,0.203) (2.571,0.212) (2.672,0.221) (2.773,0.230) (2.874,0.238) (2.976,0.246) (3.077,0.254) (3.178,0.262) (3.279,0.269) (3.380,0.277) (3.482,0.284) (3.583,0.291) (3.684,0.297) (3.785,0.304) (3.887,0.310) (3.988,0.317) (4.089,0.323) (4.190,0.329) (4.291,0.334) (4.393,0.340) (4.494,0.345) (4.595,0.351) (4.696,0.356) (4.798,0.361) (4.899,0.366) (5.000,0.371)};
  \addplot[thick,color=thmcol] coordinates {(0.040,0.000) (0.141,0.000) (0.242,0.000) (0.344,0.000) (0.445,0.000) (0.546,0.000) (0.647,0.000) (0.749,0.001) (0.850,0.001) (0.951,0.002) (1.052,0.003) (1.153,0.005) (1.255,0.007) (1.356,0.010) (1.457,0.013) (1.558,0.016) (1.660,0.020) (1.761,0.024) (1.862,0.028) (1.963,0.032) (2.064,0.037) (2.166,0.041) (2.267,0.046) (2.368,0.051) (2.469,0.056) (2.571,0.061) (2.672,0.066) (2.773,0.072) (2.874,0.077) (2.976,0.082) (3.077,0.087) (3.178,0.092) (3.279,0.098) (3.380,0.103) (3.482,0.108) (3.583,0.113) (3.684,0.118) (3.785,0.123) (3.887,0.128) (3.988,0.133) (4.089,0.138) (4.190,0.143) (4.291,0.148) (4.393,0.152) (4.494,0.157) (4.595,0.162) (4.696,0.166) (4.798,0.171) (4.899,0.175) (5.000,0.180)};
  \legend{{$x=1$},{$x=2$},{$x=3$}}
\end{axis}
\end{tikzpicture}\hfill
\begin{tikzpicture}
\begin{axis}[width=0.47\linewidth,height=5cm,xlabel={$t$},ylabel={$f(t\mid x)$},xmin=0,xmax=5,ymin=0,ymax=0.5,
  title={Density of the hitting time $\tau_x$},title style={font=\footnotesize},
  legend style={at={(0.97,0.97)},anchor=north east,font=\scriptsize},grid=major,grid style={black!10}]
  \addplot[thick,color=defcol] coordinates {(0.050,0.002) (0.134,0.195) (0.218,0.395) (0.302,0.459) (0.386,0.456) (0.469,0.428) (0.553,0.393) (0.637,0.358) (0.721,0.326) (0.805,0.297) (0.889,0.271) (0.973,0.249) (1.057,0.229) (1.141,0.211) (1.225,0.196) (1.308,0.182) (1.392,0.170) (1.476,0.159) (1.560,0.149) (1.644,0.140) (1.728,0.132) (1.812,0.124) (1.896,0.117) (1.980,0.111) (2.064,0.106) (2.147,0.100) (2.231,0.096) (2.315,0.091) (2.399,0.087) (2.483,0.083) (2.567,0.080) (2.651,0.077) (2.735,0.073) (2.819,0.071) (2.903,0.068) (2.986,0.065) (3.070,0.063) (3.154,0.061) (3.238,0.059) (3.322,0.057) (3.406,0.055) (3.490,0.053) (3.574,0.051) (3.658,0.050) (3.742,0.048) (3.825,0.047) (3.909,0.045) (3.993,0.044) (4.077,0.043) (4.161,0.042) (4.245,0.041) (4.329,0.039) (4.413,0.038) (4.497,0.037) (4.581,0.036) (4.664,0.036) (4.748,0.035) (4.832,0.034) (4.916,0.033) (5.000,0.032)};
  \addplot[thick,color=intcol] coordinates {(0.050,0.000) (0.134,0.000) (0.218,0.001) (0.302,0.006) (0.386,0.019) (0.469,0.035) (0.553,0.052) (0.637,0.068) (0.721,0.081) (0.805,0.092) (0.889,0.100) (0.973,0.106) (1.057,0.111) (1.141,0.113) (1.225,0.115) (1.308,0.116) (1.392,0.115) (1.476,0.115) (1.560,0.114) (1.644,0.112) (1.728,0.110) (1.812,0.108) (1.896,0.106) (1.980,0.104) (2.064,0.102) (2.147,0.100) (2.231,0.098) (2.315,0.095) (2.399,0.093) (2.483,0.091) (2.567,0.089) (2.651,0.087) (2.735,0.085) (2.819,0.083) (2.903,0.081) (2.986,0.079) (3.070,0.077) (3.154,0.076) (3.238,0.074) (3.322,0.072) (3.406,0.071) (3.490,0.069) (3.574,0.067) (3.658,0.066) (3.742,0.065) (3.825,0.063) (3.909,0.062) (3.993,0.061) (4.077,0.059) (4.161,0.058) (4.245,0.057) (4.329,0.056) (4.413,0.055) (4.497,0.054) (4.581,0.053) (4.664,0.052) (4.748,0.051) (4.832,0.050) (4.916,0.049) (5.000,0.048)};
  \addplot[thick,color=thmcol] coordinates {(0.050,0.000) (0.134,0.000) (0.218,0.000) (0.302,0.000) (0.386,0.000) (0.469,0.000) (0.553,0.001) (0.637,0.002) (0.721,0.004) (0.805,0.006) (0.889,0.009) (0.973,0.012) (1.057,0.016) (1.141,0.019) (1.225,0.022) (1.308,0.026) (1.392,0.029) (1.476,0.032) (1.560,0.034) (1.644,0.037) (1.728,0.039) (1.812,0.041) (1.896,0.043) (1.980,0.044) (2.064,0.046) (2.147,0.047) (2.231,0.048) (2.315,0.049) (2.399,0.049) (2.483,0.050) (2.567,0.050) (2.651,0.051) (2.735,0.051) (2.819,0.051) (2.903,0.051) (2.986,0.051) (3.070,0.051) (3.154,0.051) (3.238,0.051) (3.322,0.051) (3.406,0.051) (3.490,0.051) (3.574,0.050) (3.658,0.050) (3.742,0.050) (3.825,0.049) (3.909,0.049) (3.993,0.049) (4.077,0.048) (4.161,0.048) (4.245,0.047) (4.329,0.047) (4.413,0.047) (4.497,0.046) (4.581,0.046) (4.664,0.045) (4.748,0.045) (4.832,0.044) (4.916,0.044) (5.000,0.044)};
  \legend{{$x=1$},{$x=2$},{$x=3$}}
\end{axis}
\end{tikzpicture}
\caption{Left: $\Prob[M(t)>x]=2[1-\Phi(x/\sqrt t)]$ (file \file{lec14\_6}); at $t=5$ the values are $0.655$, $0.371$, $0.180$.
Right: the hitting-time density \eqref{ch11:eq-hitcdf} (file \file{lec14\_7}); the peak of $x=1$ is at $t=1/3$ with height
$0.463$, and levels farther away are reached later, with a much more spread-out distribution
\ts{24}{14}{31:35}\ts{24}{14}{33:44}.}\label{ch11:fig-max}
\end{figure}

\begin{intuition}[A path reaches everything, but takes an infinite mean time]
Recurrence and infinite mean go together: Brownian motion in one dimension reaches every level with probability one,
but the waiting time has the heavy tail $t^{-1/2}$. It is the continuous version of the simple random walk, which returns
to $0$ a.s.\ while the expected return time is infinite (\cref{ch04:sec-ruin}, \cref{ch04:thm-ost}: optional stopping fails for $\tau_a$ because
$B(\tau_a)=a\ne0=\E[B(0)]$ and the stopped process is not uniformly bounded).
The scaling $\tau_a\eqd a^2\tau_1$ is diffusion again: time to travel $a$ grows like $a^2$.
\end{intuition}

\subsection{Reflected and absorbed Brownian motion}\label{ch11:sec-absorbed}
Two modifications of $B$ appear in the lecture \ts{24}{14}{33:44}\ts{24}{14}{34:50}.

\begin{definition}[Reflected and absorbed Brownian motion]\label{ch11:def-refabs}
\emph{Reflected Brownian motion} at zero is $R(t)=|B(t)|$ (the path is folded onto the positive half line).
\emph{Absorbed Brownian motion} started at $x>0$ is $A(t)=x+B(t)$ for $t\le\tau$ and $A(t)=0$ for $t>\tau$, where $\tau$ is
the first hitting time of $0$.
\end{definition}

\begin{proposition}[Moments of $R$ and law of $A$]\label{ch11:prop-refabs}
\leavevmode
\begin{enumerate}[(i)]
  \item $\E[R(t)]=\sqrt{2t/\pi}$ and $\Var R(t)=(1-2/\pi)\,t$.
  \item The law of $A(t)$ has an atom at $0$ of mass $2[1-\Phi(x/\sqrt t)]$ (the probability of having been absorbed), and on $y>0$ the
      density
      \begin{equation}\label{ch11:eq-absorbed}
      p_A(y,t\mid x)=p(y,t\mid x)-p(-y,t\mid x)=\frac1{\sqrt{2\pi t}}\Bigl[e^{-(y-x)^2/2t}-e^{-(y+x)^2/2t}\Bigr].
      \end{equation}
\end{enumerate}
\end{proposition}
\begin{proof}
(i) $R(t)\eqd M(t)$ by \cref{ch11:prop-max}, so $\E[R]=\sqrt{2t/\pi}$; $\E[R^2]=\E[B^2]=t$ gives the variance.

(ii) Absorption at $0$ from $x$ is the hitting of the level $-x$ by $B$ from $0$, so $\Prob[\tau\le t]=2[1-\Phi(x/\sqrt t)]$
by \eqref{ch11:eq-hitcdf}. For $y>0$, $\Prob[x+B(t)\in\dd y,\ \tau\le t]$ is, by reflection in $0$, the density of a path ending at $-y$, that is
$p(-y,t\mid x)$; subtracting this from the free density gives \eqref{ch11:eq-absorbed} (the \emph{method of images}, with an image
charge of opposite sign at $-x$).
\end{proof}
Check: $\int_0^\infty p_A\dd y=\Phi(x/\sqrt t)-[1-\Phi(x/\sqrt t)]=1-2[1-\Phi(x/\sqrt t)]$, so density plus atom sum to $1$; and
$\E[A(t)]=x$ for all $t$ (absorbed Brownian motion is the martingale $B(t\wedge\tau)$, the loss of mass to $0$ being
compensated by the surviving paths). \emph{Not to be confused:} the slides use the word ``reflected'' both for the folded process $|B|$ and, in the section
on the reflection principle, for the path reflected \emph{in a level after the hitting time}; only the second is used in \cref{ch11:lem-reflection}.

\begin{finance}[A price that can go bankrupt]
The lecture proposes absorbed Brownian motion as the price of an asset which is positive as long as the firm is viable, but stays at $0$
once bankrupt \ts{24}{14}{34:50}. Structural credit models (Merton, Black--Cox; cf.\ \cref{ch:counterparty}) do exactly this: the
probability of default before $t$ is a hitting probability, and the survival probability with $\sigma$-volatile assets started $x$ above the barrier
is $2\Phi(x/(\sigma\sqrt t))-1$, a straightforward application of \eqref{ch11:eq-hitcdf}.
\end{finance}


% ==================================================================
\section{Brownian motion with drift and the gambler's ruin problem}\label{ch11:sec-drift}

\begin{definition}[Brownian motion with drift]\label{ch11:def-drift}
$X(t)=\mu t+\sigma B(t)$ for a standard Brownian motion $B$, a \emph{drift} $\mu\in\R$ and a \emph{volatility} $\sigma>0$
\ts{24}{14}{43:19}\ts{13}{17}{47:24}. $X$ has continuous paths, independent increments and
$X(t)-X(s)\sim N(\mu(t-s),\sigma^2(t-s))$, so $\E[X(t)]=\mu t$ and $\Var X(t)=\sigma^2t$; it reduces to $B$ for $\mu=0,\sigma=1$.
\end{definition}

\begin{proposition}[Drift wins in the long run]\label{ch11:prop-drift}
For every $\varepsilon>0$, with probability tending to one $(\mu-\varepsilon)t\le X(t)\le(\mu+\varepsilon)t$ as $t\to\infty$; indeed
$X(t)/t\to\mu$ (almost surely). Yet at moderate times the noise $\sigma\sqrt t$ dominates the drift $\mu t$ as long as $t<(\sigma/\mu)^2$.
\end{proposition}
\begin{proof}
$X(t)/t-\mu=\sigma B(t)/t$ has variance $\sigma^2/t\to0$ (convergence in probability; the almost sure version is the law of large numbers for $B$). The
two terms $\mu t$ and $\sigma\sqrt t$ are equal at $t=(\sigma/\mu)^2$.
\end{proof}
The last remark of the proposition is the reason why drifts are so hard to estimate: for a stock with $\mu=10\%$ and $\sigma=30\%$ the time
scale is $(\sigma/\mu)^2=9$ years, and for it the estimate $\hat\mu=X(T)/T$ has standard error $\sigma/\sqrt T$ (\cref{ch:volatility}).

\paragraph{One-step analysis.} Look at the increment $\Delta X=X(t+\Delta t)-X(t)=\mu\Delta t+\sigma\Delta B$ with
$\Delta B\sim N(0,\Delta t)$ independent of the past. Then
\begin{equation}\label{ch11:eq-onestep}
  \E[\Delta] X=\mu\,\Delta t,\qquad \E[(\Delta X)^2]=\sigma^2\Delta t+(\mu\Delta t)^2=\sigma^2\Delta t+o(\Delta t),\qquad
  \E[|\Delta X|^c]=o(\Delta t)\ \ (c>2)
\end{equation}
\ts{24}{14}{44:20}\ts{24}{14}{45:27}\ts{24}{14}{46:33}\ts{24}{14}{47:36}: to order $\Delta t$, only the mean and the variance of the
increment matter. If $f$ is twice differentiable, Taylor's formula and \eqref{ch11:eq-onestep} give
\begin{equation}\label{ch11:eq-generator}
  \E\bigl[f(x+\Delta X)\bigr]=f(x)+\Bigl(\mu f'(x)+\tfrac12\sigma^2f''(x)\Bigr)\Delta t+o(\Delta t),
\end{equation}
so $\mathcal Lf=\mu f'+\tfrac12\sigma^2f''$ is the \emph{generator} of the process. This expansion is why the lecture cares about
the powers of the increment: the value of a derivative on an underlying $X$ is a function $f(X)$, and its dynamics are read
from the Taylor series \ts{24}{14}{47:36}\ts{24}{14}{48:45}. It is the seed of It\^o's lemma (\cref{ch11:sec-gbm}).

\subsection{Gambler's ruin for Brownian motion}\label{ch11:sec-ruin}
Let $X$ be a Brownian motion with drift $\mu$ and volatility $\sigma$, started at $x\in(a,b)$, and let
$\tau=\min\{u:X(u)=a\text{ or }b\}$. Find
\[
  u(x)=\Prob\bigl[X(\tau)=b\mid X(0)=x\bigr],
\]
the probability of reaching the upper level first. This generalises the random walk of \cref{ch04:sec-ruin} to continuous
time and continuous state \ts{24}{14}{48:45}\ts{24}{14}{49:47}\ts{24}{14}{50:51}.

\paragraph{Infinitesimal derivation (the lecture).} Wait a short time $\Delta t$ so small that hitting a level in
the meantime is negligible. By the Markov property, $u(x)=\E[u(x+\Delta X)]$; expanding with \eqref{ch11:eq-generator},
\begin{gather*}
  u(x)=u(x)+\bigl(\mu u'(x)+\tfrac12\sigma^2u''(x)\bigr)\Delta t+o(\Delta t)\\
  \Longrightarrow\ \mu\,u'(x)+\tfrac12\sigma^2u''(x)=0,\qquad u(a)=0,\ u(b)=1
\end{gather*}
\ts{24}{14}{51:55}\ts{24}{14}{52:59}\ts{24}{14}{54:04}\ts{24}{14}{55:16}. This is a first-order equation in $v=u'$: $v'/v=-2\mu/\sigma^2$, so $u'(x)=C\,e^{-2\mu x/\sigma^2}$ and
\begin{equation}\label{ch11:eq-ruin}
  u(x)=\frac{g(x)-g(a)}{g(b)-g(a)},\quad g(x)=e^{-2\mu x/\sigma^2}\ \ (\mu\ne0);\qquad
  u(x)=\frac{x-a}{b-a}\ \ (\mu=0)
\end{equation}
\ts{24}{14}{56:16}\ts{24}{14}{57:18}\ts{24}{14}{58:18}\ts{24}{14}{59:20}. For $\mu=0$, $g(x)=x$ is the other solution of the equation
(the general solution is $Ax+B$); with no drift, the probability of success is linear in the distance travelled, exactly as for the fair gambler
(\cref{ch04:prop-ruin}).

\paragraph{Rigorous derivation by a martingale.} The infinitesimal argument assumes that $u$ is twice differentiable, which
one does not want to take for granted. The proof by optional stopping needs nothing of the kind. By
\cref{ch11:prop-elem}(iii) with $\lambda=\theta\sigma$, the process $e^{\theta X(t)}=e^{\theta\mu t+\theta\sigma B(t)}$ is a martingale iff $\theta\mu=-\theta^2\sigma^2/2$, i.e.
$\theta=-2\mu/\sigma^2$. Hence $M(t)=g(X(t))=e^{-2\mu X(t)/\sigma^2}$ is a martingale, exactly the continuous analogue of
$r^{S_n}$ with $r=q/p$ of \cref{ch04:ex-m4}. The stopped martingale $M(t\wedge\tau)$ is bounded, and $\tau<\infty$ a.s., so
optional stopping (\cref{ch04:thm-ost}(ii)) gives $g(x)=\E[g(X(\tau))]=u\,g(b)+(1-u)\,g(a)$, which is \eqref{ch11:eq-ruin}. It is also the exact limit
of the discrete formula \cref{ch04:prop-biased}: for steps $\pm\sigma\sqrt h$ with $p=\tfrac12(1+\mu\sqrt h/\sigma)$, one has
$r=q/p\approx e^{-2\mu\sqrt h/\sigma}$, and $r^{S_n}\to e^{-2\mu X/\sigma^2}$. A numerical check (not in the course): for
$\mu=0.5$, $\sigma=1$, $a=0$, $b=3$, $x=1$ formula \eqref{ch11:eq-ruin} gives $u=0.66524$, and the exact ruin probability
of the biased walk with $h=10^{-4}$ is $0.66524$.

\begin{figure}[ht]
\centering
\begin{tikzpicture}
\begin{axis}[width=0.9\linewidth,height=5.6cm,xlabel={starting level $x$ ($a=0$, $b=1$, $\sigma=1$)},ylabel={$u(x)=\Prob[\text{hit }b\text{ before }a]$},
  xmin=0,xmax=1,ymin=0,ymax=1.02,legend style={at={(0.5,1.03)},anchor=south,legend columns=5,font=\footnotesize},grid=major,grid style={black!10}]
  \addplot[thick,color=thmcol,domain=0:1,samples=60] {(exp(-2*(-3)*x)-1)/(exp(-2*(-3))-1)};
  \addplot[thick,color=fincol,domain=0:1,samples=60] {(exp(-2*(-1)*x)-1)/(exp(-2*(-1))-1)};
  \addplot[thick,color=black!70,domain=0:1,samples=2] {x};
  \addplot[thick,color=intcol,domain=0:1,samples=60] {(exp(-2*1*x)-1)/(exp(-2*1)-1)};
  \addplot[thick,color=defcol,domain=0:1,samples=60] {(exp(-2*3*x)-1)/(exp(-2*3)-1)};
  \legend{{$\mu=-3$},{$\mu=-1$},{$\mu=0$},{$\mu=1$},{$\mu=3$}}
\end{axis}
\end{tikzpicture}
\caption{Probability of hitting $b=1$ before $a=0$ as a function of the starting point, for several drifts (the four panels of
\file{lec14\_8} show $b=3$ and $b=1$ with $\mu=\pm(10^{-4},0.5,\dots,3)$). With no drift the probability is the straight line
$x$; a positive drift bends the curve above it (the barrier $b$ is reached first from almost everywhere), a negative drift below it
\ts{24}{14}{1:00:20}\ts{24}{14}{1:01:30}\ts{24}{14}{1:02:33}.}\label{ch11:fig-ruin}
\end{figure}

\begin{proposition}[Consequences of \eqref{ch11:eq-ruin}; not discussed in class]\label{ch11:prop-ruincons}
\leavevmode
\begin{enumerate}[(i)]
  \item Symmetry: $u_\mu(x)+u_{-\mu}(a+b-x)=1$.
  \item If $\mu>0$, the probability that $X$ ever falls to a level $d$ below its start is $e^{-2\mu d/\sigma^2}$,
      so the running minimum of $X$ is exponentially distributed with rate $2\mu/\sigma^2$; if $\mu\le0$ every level below is
      reached with probability one.
  \item If $\mu=0$, $\E[\tau]=(x-a)(b-x)/\sigma^2$.
\end{enumerate}
\end{proposition}
\begin{proof}
(i) reflect the interval, $x\mapsto a+b-x$, which changes the sign of the drift and swaps the two barriers.

(ii) Let $b\to\infty$ in
$1-u=(g(b)-g(x))/(g(b)-g(a))$: for $\mu>0$, $g(b)\to0$ and the limit is $g(x)/g(a)=e^{-2\mu(x-a)/\sigma^2}$; for $\mu<0$, $g(b)\to\infty$ and the limit is $1$.

(iii) $X(t)^2-\sigma^2t$ is a martingale (\cref{ch11:prop-elem}(iii)); optional stopping gives $\sigma^2\E[\tau]=\E[X(\tau)^2]-x^2
=\frac{(b-x)a^2+(x-a)b^2}{b-a}-x^2=(x-a)(b-x)$.
\end{proof}

\begin{finance}[How likely is a losing streak? (numbers not from the course)]
Take the parameters used in 2024 PS5 for a simulation, $\mu=0.30$ and $\sigma=0.40$, for the \emph{log} value of a strategy (the drift
already includes any $-\sigma^2/2$ correction). The probability that its value ever falls by $0.2$ (about $18\%$ in level terms) below its running
starting point, however long one waits, is $e^{-2\cdot0.30\cdot0.2/0.16}=e^{-0.75}=47\%$, even though the strategy has a positive expected return of
$30\%$ per year. Positive drift protects in the long run but not from drawdowns: the minimum of the path is exponential with mean $\sigma^2/(2\mu)=0.27$.
\end{finance}

% ==================================================================
\section{Geometric Brownian motion and the first look at It\^o's lemma}\label{ch11:sec-gbm}

The price of a stock cannot be $S_0+\sigma B_t$: it would become negative. What is reasonable, the lecture argues, is that the
\emph{percentage} change is normal, $\Delta S/S=\mu\Delta t+\sigma\Delta B$, written $\dd S_t=S_t(\mu\dd t+\sigma\dd B_t)$
\ts{13}{17}{1:00:19}\ts{13}{17}{1:01:24}. The natural question, asked in the 2013 lecture, is whether the solution of $\dd S=S\dd B$ is $S_t=e^{B_t}$, as in classical calculus; the answer is no
\ts{13}{17}{1:01:24}\ts{13}{17}{1:02:28}. Here is why.

\paragraph{The Taylor argument.} Let $f$ be a smooth function (in the 2013 lecture a call payoff, $f(B_T)=(B_T-s_0)^+$, or the price of any derivative on the underlying)
and $B$ a Brownian motion \ts{13}{17}{1:03:29}\ts{13}{17}{1:04:29}\ts{13}{17}{1:05:33}. If $B$ were differentiable, $\dd f(B_t)=f'(B_t)\frac{\dd B}{\dd t}\dd t$; it is not
\ts{13}{17}{1:06:36}, so one keeps the increment $\dd B$ and writes Taylor's formula
$f(B_{t+x})-f(B_t)=f'(B_t)\,\Delta B+\tfrac12f''(B_t)\,\Delta B^2+\dots$ with $\Delta B=B_{t+x}-B_t$
\ts{13}{17}{1:07:39}\ts{13}{17}{1:08:44}\ts{13}{17}{1:09:46}\ts{13}{17}{1:10:48}. In calculus the term $\Delta B^2$ is negligible; here, by
\eqref{ch11:eq-dB2}, $\Delta B^2\approx\Delta t$ is of the same order as the time step, so it cannot be dropped:
\begin{equation}\label{ch11:eq-ito}
  \dd f(B_t)=f'(B_t)\dd B_t+\tfrac12f''(B_t)\dd t\qquad\text{(It\^o's lemma, heuristic form)}
\end{equation}
\ts{13}{17}{1:11:55}\ts{13}{17}{1:13:01}\ts{13}{17}{1:15:12}. The lecturer's summary: if one thing is to be remembered it is \eqref{ch11:eq-ito}, if two, that $(\dd B)^2=\dd t$
\ts{13}{17}{1:14:05}. Proofs and generalisations (functions of $t$ and $X$, stochastic integrals) are in \cref{ch:stochcalc}.

\paragraph{Application to the exponential.} With $f(x)=e^{x}$, \eqref{ch11:eq-ito} gives $\dd e^{B_t}=e^{B_t}\dd B_t+\tfrac12e^{B_t}\dd t$: the
process $e^{B}$ has a \emph{drift} $\tfrac12S\dd t$ that $\dd S=S\dd B$ does not have. Indeed $\E[e^{B_t}]=e^{t/2}$ (moment generating function of the normal), whereas
$\dd S=S\dd B$ means that $S$ is a martingale with constant mean $S_0$; the same computation is the exponential martingale of
\cref{ch11:prop-elem}(iii), and the correction $-\tfrac12\sigma^2t$ in the exponent is the same lognormal correction as in \cref{ch04:rem-logprice}.
The solution of $\dd S=S(\mu\dd t+\sigma\dd B)$ is therefore
\begin{equation}\label{ch11:eq-gbm}
  S_t=S_0\exp\Bigl[\bigl(\mu-\tfrac12\sigma^2\bigr)t+\sigma B_t\Bigr],
\end{equation}
\emph{geometric Brownian motion} (GBM): $\log S_t$ is a Brownian motion with drift $\mu-\tfrac12\sigma^2$ and volatility $\sigma$, so $S_t$ is lognormal
(\cref{ch03:sec-normal}) with
\[
  \E[S_t]=S_0e^{\mu t},\qquad \operatorname{median}(S_t)=S_0e^{(\mu-\sigma^2/2)t},\qquad \Var S_t=S_0^2e^{2\mu t}\bigl(e^{\sigma^2t}-1\bigr).
\]

\paragraph{The mechanism in discrete time (our check).} The discrete product of one-step relative changes,
$\prod_i(1+\sigma\Delta_iB)$ over a grid of mesh $h$, has mean exactly $\prod(1+0)=1$ and equals
$\exp\bigl(\sum_i\ln(1+\sigma\Delta_iB)\bigr)=\exp\bigl(\sigma B_t-\tfrac12\sigma^2\sum_i(\Delta_iB)^2+\dots\bigr)$; the sum of
squares converges to $t$ by \cref{ch11:thm-qv}, so the product converges to $\exp(\sigma B_t-\tfrac12\sigma^2t)$: the extra
factor $e^{-\sigma^2t/2}$ is the quadratic variation at work. In a simulation with $2\times10^5$ paths of $400$ $\pm\sqrt{h}$ steps ($t=1$, $\sigma=1$) the
mean of the product is $1.005$, whereas the mean of $e^{B_1}$ is $1.656$ (theory: $1$ and $e^{1/2}=1.649$).

\begin{finance}[Arithmetic versus geometric growth]
With $\mu=8\%$ and $\sigma=20\%$ the expected value of an investment grows as $e^{0.08\cdot10}=2.23$ over ten years, but the median as
$e^{(0.08-0.02)10}=1.82$, and the probability of ending below the initial value is $\Phi(-0.06\cdot\sqrt{10}/0.2)=17\%$. The gap between the mean and the typical
path grows with $\sigma^2t$, exactly as for the multiplicative martingales of \cref{ch04:sec-mart} (volatility drag). GBM is the model of Black, Scholes and Merton
(\cref{ch:bs}); Brownian motion with drift, on the other hand, is the model of the \emph{logarithm} of the price, which is why the
gambler's-ruin formulas of \cref{ch11:sec-ruin} apply to log-prices and log-levels.
\end{finance}

% ==================================================================
\section{Processes built from Brownian motion}\label{ch11:sec-variants}

\subsection{The Brownian bridge and the empirical distribution function}\label{ch11:sec-bridge}
\begin{definition}[Brownian bridge]\label{ch11:def-bridge}
$X(t)=B(t)-tB(1)$, $0\le t\le1$ \ts{24}{14}{35:57}\ts{24}{14}{37:05}.
\end{definition}
The path is pinned at both ends: $X(0)=0=X(1)$, since it is $B$ minus the straight line from $0$ to $B(1)$ that it would have followed on average. It is a centred Gaussian
process with
\begin{equation}\label{ch11:eq-bridge}
  \Var X(t)=t(1-t),\qquad \Cov(X(s),X(t))=\min(s,t)-st\quad(s,t\in[0,1]),
\end{equation}
from $\Cov(X(s),X(t))=s-st-st+st$ for $s\le t$. Since $\Cov(X(t),B(1))=t-t=0$ and the pair is jointly Gaussian, $X$ is independent of $B(1)$;
consequently $X$ has the law of Brownian motion \emph{conditioned} to return to $0$ at time $1$ (the endpoint is
irrelevant for the shape of the fluctuations). The variance is largest in the middle, $\tfrac14$ at $t=\tfrac12$.

\paragraph{Where it comes from: the empirical distribution function.} Let $U_1,\dots,U_n$ be i.i.d.\ uniform on $(0,1)$ and
$\hat F_n(t)=\#\{U_i\le t\}/n$. Each indicator $\ind\{U_i\le t\}$ is Bernoulli with mean $t$ and variance $t(1-t)$, so $\E[\hat F_n(t)]=t$ and
$\Var\hat F_n(t)=t(1-t)/n$; by the central limit theorem
$\sqrt n[\hat F_n(t)-t]\to N(0,t(1-t))$ \ts{24}{14}{38:05}\ts{24}{14}{39:08}\ts{24}{14}{40:09}. The whole process $\sqrt n[\hat F_n(t)-t]$ converges to
the Brownian bridge (Donsker, 1952): variance $t(1-t)$ and covariance $\min(s,t)-st$ are those of \eqref{ch11:eq-bridge}, checked directly
on the multinomial counts.

\begin{finance}[Testing a return distribution: probability integral transform and Kolmogorov--Smirnov]
If $Y_1,\dots,Y_n$ are i.i.d.\ with a \emph{continuous} cdf $F$, then $U_i=F(Y_i)$ are i.i.d.\ uniform (\cref{ch03:prop-pit}), so
the statistic $D_n=\sup_y|\hat F_n(y)-F(y)|=\sup_t|\hat F_n^U(t)-t|$ has a distribution that does not depend on $F$: $\sqrt nD_n\to\sup_{t}|X(t)|$, whose law is
$\Prob[\sup|X|\le x]=1-2\sum_{k\ge1}(-1)^{k-1}e^{-2k^2x^2}$ (Kolmogorov); the $5\%$ critical value is $1.358$.
To test whether a set of returns is compatible with an assumed model, transform the data by the model's cdf and test uniformity
\ts{24}{14}{41:10}\ts{24}{14}{42:14}\ts{24}{14}{43:19}. (If the parameters of $F$ are estimated from the same data, the critical values are smaller, as in the Lilliefors
correction.) This is a model-selection tool; residual diagnostics of \cref{ch:timeseries} play the same role for time series.
\end{finance}

\subsection{Sampling at random times: Laplace increments}\label{ch11:sec-laplace}
Observe $B$ at $t=0,1,2,\dots$: the increments are i.i.d.\ $N(0,1)$ (with drift and volatility, $N(\mu,\sigma^2)$)
\ts{24}{14}{1:13:28}. Suppose instead the observation times are random, with i.i.d.\ exponential gaps $T\sim\mathrm{Exp}(\lambda)$ (rate $\lambda$)
independent of $B$ \ts{24}{14}{1:14:35}. Then the increment $B(T)$ has density
\begin{equation}\label{ch11:eq-laplace}
  f(x)=\int_0^\infty\lambda e^{-\lambda t}\frac{e^{-x^2/2t}}{\sqrt{2\pi t}}\dd t=\sqrt{\tfrac\lambda2}\;e^{-\sqrt{2\lambda}\,|x|},
\end{equation}
the \emph{Laplace} (two-sided exponential) density, instead of a bell curve: a \emph{normal variance mixture} in
the sense of \cref{ch03:prop-mixture}, $B(T)=\sqrt T\,Z$, with variance $\E[T]=1/\lambda$ and excess kurtosis $3\E[T^2]/(\E[T])^2-3=3$ (checked by numerical integration of \eqref{ch11:eq-laplace} at two points). Its tails are heavier than the normal ones, and the
lecture reports that a Laplace model fits S\&P~500 returns much better \ts{24}{14}{1:15:44}. The economic reading: the useful clock for prices
is not calendar time but a random \emph{business time} (trades, information arrivals), and the process is a Brownian motion \emph{subordinated} to a
random clock. Time-changed and stochastic-volatility models take this idea further (\cref{ch:volatility}).

\subsection{Two-dimensional Brownian motion}\label{ch11:sec-2d}
The lecture opens with the animation of a particle moving in the plane, where each coordinate takes independent Gaussian steps,
no direction being preferred \ts{24}{14}{1:04}\ts{24}{14}{2:04}\ts{24}{14}{3:08}. Formally $\bm B(t)=(B_1(t),B_2(t))$ with independent standard Brownian motions has density
$(2\pi t)^{-1}e^{-|\bm y|^2/2t}$, invariant under rotations; the distance $|\bm B(t)|$ has the Rayleigh law
$\Prob[|\bm B(t)|<\varepsilon]=1-e^{-\varepsilon^2/(2t)}$ (for $t=\varepsilon=1$: $39.3\%$) and $\E[|\bm B(t)|]=\sqrt{\pi t/2}$.
The path in the plane is fractal-looking and its two projections are the one-dimensional processes of the rest of the chapter (file \file{lec14\_2}, $2000$ steps).
Correlated Brownian motions, obtained from independent ones by a matrix (\cref{ch:linalg}), are the basis of models with several risk factors.

\subsection{Filtrations and adapted processes}\label{ch11:sec-adapted}
Let $X$ be a stochastic process, for instance a price. The \emph{filtration} $\{\mathcal F_t\}$ generated by $X$ is the increasing family of $\sigma$-fields of
events determined by $\{X(u),\,0\le u\le t\}$: $\mathcal F_t\subset\mathcal F_{t'}$ for $t<t'$; it is the information available at time $t$
\ts{24}{14}{1:16:44}\ts{24}{14}{1:17:48}. A process $Y$ is \emph{adapted} to $X$ if $Y(t)$ depends only on $\{X(u),u\le t\}$, equivalently
$\mathcal G_t\subset\mathcal F_t$ for the filtration $\mathcal G$ of $Y$: the transformation is not forward-looking. For example if $Y(t)$
is the cash flow of a strategy that trades the security $X$, adaptedness says that one cannot trade on tomorrow's price
\ts{24}{14}{1:17:48}\ts{24}{14}{1:18:55}. Adaptedness is what makes the gains from trading a martingale when $X$ is one: the continuous-time
analogue of the martingale transform of \cref{ch04:thm-transform} (with constant-mean processes, martingales, ``applied again and again'' in the
lecture \ts{24}{14}{1:18:55}\ts{24}{14}{1:19:59}) is the stochastic integral $\int Y\dd X$ of \cref{ch:stochcalc}, and the price of an option is a conditional
expectation with respect to $\mathcal F_t$ (\cref{ch:bs}). All the martingale statements of \cref{ch11:prop-elem} hold with respect to the filtration of $B$.

\subsection{L\'evy processes: Brownian motion and the Poisson process}\label{ch11:sec-levy}
2024 PS5 opens with a general definition (Problems~1--2, \cref{ch11:ex-levy-bm,ch11:ex-levy-poisson}); the topic was not lectured on.

\begin{definition}[L\'evy process]\label{ch11:def-levy}
$\{L(t),\,t\ge0\}$ is a \emph{L\'evy process} if
\begin{enumerate}[(i)]
  \item $L(0)=0$;
  \item stationary increments: $L(t)-L(s)\eqd L(t-s)$ for $s\le t$;
  \item independent increments over disjoint intervals;
  \item continuity in probability: $\lim_{s\to t}\Prob[|L(t)-L(s)|>\varepsilon]=0$ for all $\varepsilon>0$ and $t$.
\end{enumerate}
\end{definition}
(The problem set indexes $t$ over all of $\R$; only $t\ge0$ is used.) It generalises Definition~\ref{ch11:def-bm}, which asks in
addition that the increments be normal and the paths continuous: continuity in probability is much weaker (it does not forbid jumps
at unpredictable times). Two examples are the subject of the problem set. \emph{Brownian motion with drift} (\cref{ch11:def-drift})
is a L\'evy process: its increments are $N(\mu h,\sigma^2h)$. The \emph{Poisson process} $N$ with intensity $\lambda$ has
independent increments with $N(t)-N(s)\sim\mathrm{Poisson}(\lambda(t-s))$; it is continuous in probability, since
$\Prob[N(t)-N(s)\ge1]=1-e^{-\lambda(t-s)}\to0$, although \emph{every} path is a pure-jump staircase.
The two are the extreme cases of the L\'evy--It\^o decomposition (not covered): every L\'evy process is the sum of a drift, a Brownian motion,
and a superposition of jumps (compensated when small).

\begin{table}[ht]
\centering\small
\begin{tabularx}{\linewidth}{@{}l L L L@{}}
\toprule
& \textbf{paths} & \textbf{increment over $h$} & \textbf{quadratic variation on $[0,T]$}\\
\midrule
$X=\mu t+\sigma B$ & continuous, nowhere differentiable & $N(\mu h,\sigma^2h)$ & $\sigma^2T$, deterministic\\
Poisson $N$, rate $\lambda$ & piecewise constant, unit jumps & $\mathrm{Poisson}(\lambda h)$ & $N(T)$ (each jump adds $1^2$), random\\
\bottomrule
\end{tabularx}
\caption{Brownian motion with drift and the Poisson process, the two building blocks of jump--diffusion models.}\label{ch11:tab-levy}
\end{table}

\begin{intuition}[Two universal limits]
Sums of many small independent increments give Brownian motion (\cref{ch11:thm-donsker}). Sums of many \emph{rare} unit
increments give the Poisson process: a walk that steps up by $1$ with probability $\lambda/n$ at each of $n$ sub-intervals has
$\mathrm{Bin}(n,\lambda/n)\to\mathrm{Poisson}(\lambda)$ increments (the law of small numbers). The compensated process
$N(t)-\lambda t$ is a martingale with variance $\lambda t$, and $(N(t)-\lambda t)/\sqrt\lambda$ looks Brownian when the rate is large, a bridge between the two limits. In finance the two blocks together give the
\emph{jump--diffusion} of Merton (1976): Brownian noise for the day-to-day movements, Poisson jumps for crashes, defaults and announcements.
Realised variance (\cref{ch11:sec-paths}) measures both: the quadratic variation of the diffusive part plus the sum of the squared jumps.
\end{intuition}

% ==================================================================
\section{Looking ahead and the two lectures compared}\label{ch11:sec-ahead}
\begin{itemize}
  \item $(\dd B)^2=\dd t$ and \eqref{ch11:eq-ito} are the beginning of \emph{stochastic calculus}: It\^o integrals $\int\theta\dd B$ for adapted $\theta$, It\^o's lemma for
        $f(t,X_t)$, geometric Brownian motion as the solution of an SDE (\cref{ch:stochcalc,ch:sde}), Girsanov's theorem and risk-neutral pricing (\cref{ch:bs}).
  \item The martingale $e^{\lambda B-\lambda^2t/2}$ and the gambler's-ruin martingale $e^{-2\mu X/\sigma^2}$ are the continuous-time versions
        of \cref{ch04:ex-m4}; the reflection principle and hitting-time density are the working tools for barrier options and
        structural credit models.
  \item The diffusion equation \eqref{ch11:eq-heat} is the seed of the Black--Scholes PDE, the Fokker--Planck equation and the Feynman--Kac formula of the next chapters.
  \item The Dickey--Fuller distribution of \cref{ch09:sec-df} is the law of $\int B\dd B/\int B^2\dd t$.
\end{itemize}

\begin{coursenote}[2013 and 2024 compared; problem-set overlap]
Both years cover the same core: definition, random-walk limit, non-differentiability, reflection principle and the maximum, quadratic variation, drift. They differ
in emphasis. \emph{2013} (C.~Lee) is rigorous about the definition (existence theorem, path space), proves non-differentiability by the maximum, and spends the second half of the lecture on
\emph{quadratic variation as the reason for It\^o's lemma}, ending with the question $S_t=e^{B_t}$?; he does not treat the bridge,
Laplace increments, reflected and absorbed Brownian motion, the hitting-time density, or gambler's ruin with drift. \emph{2024}
(P.~Kempthorne) is more applied and computational: transition density and diffusion equation, the normalised random walk simulated
for Bernoulli and Gaussian steps, path pairing for the reflection principle, hitting-time density,
the bridge and the empirical distribution function, gambler's ruin by infinitesimal analysis, quadratic variation of log-prices, subordination and adapted processes; it
stops before It\^o's lemma, which is in the next lecture (\cref{ch:stochcalc}). The quadratic-variation theorem is stated for arbitrary partitions in 2024 and proved for
equal spacing in 2013; \cref{ch11:thm-qv} above proves it in general. Problem sets: 2024 PS4 Problem~1 is 2013 PS8 Problem~A-2 (same three parts, with $\sigma=2$ in part~(b)
in 2024); 2013 PS8 Problem~B-1 (density and rotation invariance of a two-dimensional Brownian motion, and $\Prob[\bm B(t)\in D]$ for a disc $D$) belongs to \cref{ch:stochcalc}, together with the rest of
PS8; \cref{ch11:sec-2d} gives the results it needs. 2024 PS4 Problems~3--4 (MLE for the geometric Brownian motion) are treated in \cref{ch:volatility}.
\end{coursenote}

% ==================================================================
\begin{keyformulas}
\begin{align*}
&\text{Brownian motion:} && B(0)=0,\ B(t)-B(s)\sim N(0,t-s)\text{ independent}\\
& && \Cov(B(s),B(t))=\min(s,t)\\
&\text{Martingales:} && B(t),\quad B(t)^2-t,\quad e^{\lambda B(t)-\lambda^2t/2}\\
&\text{Density:} && p(y,\tau\mid x)=\tfrac{1}{\sqrt{2\pi\sigma^2\tau}}e^{-(y-x)^2/(2\sigma^2\tau)}\\
& && \partial_\tau p=\tfrac{\sigma^2}{2}\partial_y^2p\\
&\text{Donsker:} && B_n(t)=S_{\lfloor nt\rfloor}/\sqrt n\ \Rightarrow\ B(t)\quad(\E[X]=0,\ \Var X=1)\\
&\text{Scaling:} && \{B(ct)\}\eqd\{\sqrt c\,B(t)\}\\
&\text{Paths:} && \text{continuous, nowhere differentiable; }\ \sum(\Delta B)^2\to T,\ \ (\dd B)^2=\dd t\\
&\text{Maximum:} && \Prob[M(t)\ge a]=2\Prob[B(t)>a]=2[1-\Phi(a/\sqrt t)]\\
&\text{Hitting time:} && f_{\tau_a}(t)=\tfrac{a}{\sqrt{2\pi}}t^{-3/2}e^{-a^2/(2t)},\quad \E[\tau_a]=\infty\\
&\text{Reflected:} && \E[|B(t)|]=\sqrt{2t/\pi},\quad \Var|B(t)|=(1-\tfrac2\pi)t\\
&\text{Absorbed at 0:} && p_A(y)=p(y,t\mid x)-p(-y,t\mid x)\\
& && \Prob[A(t)=0]=2[1-\Phi(x/\sqrt t)]\\
&\text{Bridge:} && X(t)=B(t)-tB(1),\quad \Cov(X(s),X(t))=\min(s,t)-st\\
&\text{Drift:} && X=\mu t+\sigma B,\quad \mathcal Lf=\mu f'+\tfrac12\sigma^2f''\\
&\text{Ruin:} && \Prob[X(\tau)=b\mid X(0)=x]=\tfrac{e^{-2\mu x/\sigma^2}-e^{-2\mu a/\sigma^2}}{e^{-2\mu b/\sigma^2}-e^{-2\mu a/\sigma^2}}\\
& && \bigl(=\tfrac{x-a}{b-a}\text{ if }\mu=0\bigr)\\
&\text{It\^o (heuristic):} && \dd f(B)=f'(B)\dd B+\tfrac12f''(B)\dd t\\
&\text{GBM:} && S_t=S_0e^{(\mu-\sigma^2/2)t+\sigma B_t},\quad \E[S_t]=S_0e^{\mu t}\\
&\text{Laplace:} && B(T),\ T\sim\mathrm{Exp}(\lambda):\ f(x)=\sqrt{\lambda/2}\,e^{-\sqrt{2\lambda}|x|}
\end{align*}
\end{keyformulas}

% ==================================================================
\section*{Exercises}
\addcontentsline{toc}{section}{Exercises}

\subsection*{From the problem sets}
\begin{exercise}[PS4 (2024), Problem 1; PS8 (2013), Problem A-2: conditional moments]\label{ch11:ex-cond}
\leavevmode
\begin{enumerate}[(a)]
  \item Let $B(t)$ be a standard Brownian motion and $t>s\ge0$ fixed. Compute $\E[B(t)\mid B(s)]$ and $\Var[B(t)\mid B(s)]$.
  \item Let $X(t)$ be a Brownian motion with drift $\mu$ and volatility $\sigma=2$ (the 2013 version has no volatility). Compute $\E[X(t)\mid X(s)]$ and $\Var[X(t)\mid X(s)]$.
  \item For a standard Brownian motion and fixed reals $t>s>0$, compute $\E\bigl[\exp\bigl(\sigma(B(t)-B(s))\bigr)\bigr]$, for a real constant $\sigma$ (the multiplier of the increment inside the exponential).
\end{enumerate}
\end{exercise}

\begin{exercise}[PS4 (2024), Problem 2: Brownian scaling]\label{ch11:ex-scaling}
Let $B$ be a standard Brownian motion. Prove:
\begin{enumerate}[(a)]
  \item if $B(0)=0$, then for every $t>0$, $\{B(st),s\ge0\}\eqd\{t^{1/2}B(s),s\ge0\}$ as processes;
  \item more precisely, the two families have the same finite-dimensional distributions: for $s_1<\dots<s_n$,
      $(B(s_1t),\dots,B(s_nt))\eqd(t^{1/2}B(s_1),\dots,t^{1/2}B(s_n))$.
\end{enumerate}
\end{exercise}

\begin{exercise}[PS5 (2024), Problem 1: Brownian motion with drift as a L\'evy process]\label{ch11:ex-levy-bm}
Let $W$ be a Brownian motion with drift $\mu\in\R$ and volatility $\sigma>0$.
\begin{enumerate}[(a)]
  \item Prove that $W$ is a L\'evy process in the sense of \cref{ch11:def-levy}.
  \item Simulate the path of $W$ with $T=1$ year, $\mu=0.30$, $\sigma=0.40$ and $m=252$ increments, $t_i=iT/m$, and plot five simulated paths (the process values at the time increments).
\end{enumerate}
\end{exercise}

\begin{exercise}[PS5 (2024), Problem 2: the Poisson process]\label{ch11:ex-levy-poisson}
A stochastic process $\{N(t),t\ge0\}$ is a Poisson process of intensity (jump rate) $\lambda$ if for any $0=t_0<t_1<\dots<t_n$ the increments $N(t_k)-N(t_{k-1})$ are independent,
and $N(t)-N(s)$ is Poisson with mean $\lambda(t-s)$ for all $0\le s<t$.
\begin{enumerate}[(a)]
  \item Prove that $N$ is a L\'evy process.
  \item Simulate the path of a Poisson process with $T=1$ year, $\lambda=24$ events per year and $m=252$ increments, $t_i=iT/m$, and plot five simulated paths.
\end{enumerate}
\end{exercise}

\subsection*{Additional exercises (not from the course)}
\begin{exercise}[not from the course: time inversion]\label{ch11:ex-inversion}
Prove that $W(t)=tB(1/t)$ ($t>0$), $W(0)=0$, is a standard Brownian motion (check that it is centred Gaussian with covariance $\min(s,t)$; for continuity at $0$ use
$B(u)/u\to0$ as $u\to\infty$). Deduce that $B(t)/t\to0$ as $t\to\infty$ from the continuity of $W$ at $0$.
\end{exercise}

\begin{exercise}[not from the course: reflection and the minimum]\label{ch11:ex-refl}
Let $m(t)=\min_{u\le t}B(u)$. Show $\Prob[m(t)\le-a]=2\Prob[B(t)<-a]$ and that $M(t)-B(t)\eqd M(t)\eqd|B(t)|$ at a fixed time $t$ (hint: use \cref{ch11:rem-joint} or the
time-reversed path $B(t-u)-B(t)$). What is the probability that a driftless log price with $\sigma=25\%$ never falls below $-10\%$ during one year?
\end{exercise}

\begin{exercise}[not from the course: a hitting time with infinite mean]\label{ch11:ex-tau}
Show $\Prob[\tau_a<\infty]=1$ and $\E[\tau_a]=\infty$ from \eqref{ch11:eq-hitcdf}. Why does optional stopping ($\E[B(\tau_a)]=0$)
fail? Use $\tau_a\eqd a^2\tau_1$ to show that the median of $\tau_a$ scales like $a^2$; compute the median of $\tau_1$.
\end{exercise}

\begin{exercise}[not from the course: absorbed Brownian motion]\label{ch11:ex-abs}
For absorbed Brownian motion $A$ started at $x>0$ (\cref{ch11:def-refabs}) show that \eqref{ch11:eq-absorbed} and the mass $2[1-\Phi(x/\sqrt t)]$ at $0$ sum to $1$, and that $\E[A(t)]=x$.
Is $A(t)^2-t\wedge\tau$ a martingale?
\end{exercise}

\begin{exercise}[not from the course: exit time and ruin]\label{ch11:ex-exit}
For $X=\mu t+\sigma B$ with $\mu\ne0$ compute $\E[\tau]$ for the exit time of $(a,b)$ (hint: the martingale $X(t)-\mu t$ and \eqref{ch11:eq-ruin}). Check that the limit $\mu\to0$ is $(x-a)(b-x)/\sigma^2$.
\end{exercise}

\begin{exercise}[not from the course: geometric Brownian motion]\label{ch11:ex-gbm}
For $S_t$ in \eqref{ch11:eq-gbm} compute $\E[S_t]$, $\Var S_t$, and $\Prob[S_t<S_0]$. With $\mu=0.08$, $\sigma=0.2$ find the horizon at which the probability of a loss falls below $5\%$.
Show that $\E[e^{\lambda B(t)}]=e^{\lambda^2t/2}$ and use it to explain why $e^{B}$ is not a martingale.
\end{exercise}

\begin{exercise}[not from the course: the Laplace mixture]\label{ch11:ex-laplace}
Derive \eqref{ch11:eq-laplace} (hint: $\int_0^\infty t^{-1/2}e^{-\lambda t-x^2/2t}\dd t=\sqrt{\pi/\lambda}\,e^{-\sqrt{2\lambda}|x|}$), verify that the density integrates to one, and compute its variance and excess kurtosis.
\end{exercise}
